\documentclass[12pt]{amsart}
\usepackage[a4paper,hmargin=2.5cm,vmargin=2.4cm]{geometry}
\usepackage{amsmath,amssymb,amsthm,mathtools}
\usepackage{microtype}
\usepackage[hidelinks]{hyperref}
\usepackage{enumitem}
\setlength{\emergencystretch}{2em}

\newcommand{\E}{\mathbb E}
\newcommand{\N}{\mathbb N}
\newcommand{\Z}{\mathbb Z}
\newcommand{\R}{\mathbb R}
\newcommand{\C}{\mathbb C}
\newcommand{\1}{\mathbf 1}
\newcommand{\e}{\mathrm e}
\newcommand{\T}{\mathbb T}
\newcommand{\Cramer}{\textup{Cr\'amer}}
\DeclareMathOperator{\supp}{supp}

\newtheorem{theorem}{Theorem}[section]
\newtheorem{proposition}[theorem]{Proposition}
\newtheorem{lemma}[theorem]{Lemma}
\newtheorem{corollary}[theorem]{Corollary}
\newtheorem{remark}[theorem]{Remark}
\theoremstyle{definition}
\newtheorem{definition}[theorem]{Definition}

\setcounter{tocdepth}{1}
\title{Quantitative polynomial Szemer\'edi theorems along shifted
primes with nil-Bohr restrictions}

\author{}
\date{}

\begin{document}
\begin{abstract}
For every fixed \(k\geq 2\) and every \(\beta\) in a dense subset
of \((0,1)\) consisting of Liouville numbers, we prove that, for
some \(c>0\) and all sufficiently large \(N\), every subset of
\([N]\) of density at least \((\log N)^{-c}\) contains a polynomial
configuration 
$$x,x+(p-1),x+(p-1)^2,\ldots,x+(p-1)^k$$ 
whose prime \(p\)
belongs to  the return set
\(
 \left\{p\in\mathbb P:
 \left\{\sqrt{3}\,p\lfloor\beta p\rfloor\right\}
 \in (1/4,1/2)\right\}.
\)
The corresponding multiple ergodic averages over the selected
primes converge almost everywhere to the same limit as the
unrestricted prime averages. The main contribution is a local
generalized von Neumann theorem for distinct-degree polynomial
configurations. 
\end{abstract}
\maketitle
\pagestyle{plain}

\tableofcontents

\section{Introduction}
\label{s:introduction}

Szemer\'edi's theorem states that a set of integers of positive upper
density contains arithmetic progressions of every finite length
\cite{Szemeredi}. Its polynomial extension, due to Bergelson and Leibman
\cite{BergelsonLeibman}, replaces the linear shifts by integer polynomials
with zero constant term. A further question is how much freedom remains
in the parameter producing the configuration. Can it be required to be
one less than a prime, a return time of a rotation, or a return time of a
polynomial orbit? We study configurations whose parameter is one less
than a prime and a return time of a Heisenberg nilrotation. We prove a
quantitative recurrence theorem and a pointwise comparison with the
unrestricted prime averages. Both arguments use cancellation of the
centered selector against ordinary polynomial phases.

The discrete formulation is useful for separating the two requirements.
Given $A\subset[N]:=\{1,\ldots,N\}$, we seek integers $x$ and
$n\geq1$ such that
\begin{equation}\label{e:intro-pattern}
 x,\quad x+P_1(n),\quad\ldots,\quad x+P_\ell(n)\in A.
\end{equation}
% We start with an arbitrary set $A\subset[N]:=\{1,\ldots,N\}$ and seek
% \begin{equation}\label{e:intro-pattern}
%  x,\quad x+P_1(n),\quad\ldots,\quad x+P_\ell(n)\in A.
% \end{equation}
Throughout the introduction the polynomials are fixed and satisfy
\begin{equation}\label{e:fixed-polynomials}
 P_0=0,\qquad P_i\in\mathbb Z[y],\qquad P_i(0)=0,\qquad
 1\leq\deg P_1<\cdots<\deg P_\ell=d.
\end{equation}
This leads to a the following question: what restrictions can one impose on $n$
while still guaranteeing such configurations in every sufficiently dense
set $A$? The density of the allowed
parameters alone does not settle this question, since $A$ may
itself impose arithmetic or dynamical restrictions on the
parameters that can occur. For example, the primes do not form
a set of recurrence.

We take $n=p-1$ with $p$ prime. For two-point configurations this
restriction goes back to S\'ark\"ozy \cite{Sarkozy}. Polynomial
multiple recurrence along shifted primes was established by Wooley
and Ziegler \cite{WooleyZiegler} and, in a multidimensional setting,
by Frantzikinakis, Host and Kra \cite{FHK}. We ask whether recurrence remains under a further dynamical
restriction on $p-1$.


The simplest restriction is to require $p$ to lie in a fixed arithmetic
progression, for example $p\equiv1\pmod q$. This can be expressed as
a Bohr condition: for $\alpha=1/q$ and $0<\varepsilon<1/q$,
\[
 \|\alpha(p-1)\|_{\mathbb T}<\varepsilon
 \quad\Longleftrightarrow\quad p\equiv1\pmod q.
\]
Allowing arbitrary real frequencies generalizes these congruence
conditions to simultaneous Bohr conditions
\begin{align}\label{e:bohrcondition}
 \|\alpha_j(p-1)\|_{\mathbb T}<\varepsilon
 \qquad(1\leq j\leq r).
\end{align}
Dynamically, these require the rotation by
$(\alpha_1,\ldots,\alpha_r)$ on $\mathbb T^r$ to return close to its
starting point at time $p-1$.
When the configuration contains a linear shift, an elementary pigeonhole
argument allows us to force these Bohr conditions and deduce the result
from the shifted-prime theorem.

\begin{corollary}[Bohr restrictions]\label{c:bohr}
 Assume \eqref{e:fixed-polynomials} and $P_1(y)=ay$ with
 $a\in\mathbb Z\setminus\{0\}$. Fix real numbers
 $\alpha_1,\ldots,\alpha_r$ and $\varepsilon>0$. There exists
 $\kappa>0$ such that, for every sufficiently large $N$, every
 $A\subset[N]$ with $|A|\geq N(\log N)^{-\kappa}$ contains a configuration \eqref{e:intro-pattern} of pairwise distinct
points, with $n=p-1$ for a prime $p$ satisfying \eqref{e:bohrcondition}.
\end{corollary}

\begin{proof}
Let $B_1,\ldots,B_K$ partition $\mathbb T^r$ into half-open boxes
of side length less than $\varepsilon$, and put
\[
 A_i:=\left\{x\in A:
 \left(\frac{\alpha_1x}{a},\ldots,\frac{\alpha_rx}{a}\right)
 \pmod1\in B_i\right\}.
\]
Choose a largest $A_i$, so that $|A_i|\ge |A|/K$. Since $K$ is
independent of $N$, choosing $\kappa$ smaller than the exponent in
\cite[Theorem~1.2]{KMTT} allows us to apply that theorem to $A_i$
for sufficiently large $N$.

The theorem gives a configuration of pairwise distinct points
inside $A_i$ with parameter $p-1$ for some prime $p$.
In particular, the images of $x$ and $x+a(p-1)$ lie in the same box
$B_i$. Their $j$th coordinates differ by $\alpha_j(p-1)$ modulo one,
and hence
\[
 \|\alpha_j(p-1)\|_{\mathbb T}<\varepsilon
 \qquad(1\le j\le r).
\]
\end{proof}

This argument works because subtracting the linear torus
coordinates eliminates the starting point $x$. A natural attempt to
force a quadratic condition is to partition $A$ according to
$\alpha x^2\pmod1$. If $x$ and $x+n$ lie in the same cell, we obtain
\[
 \|\alpha(x+n)^2-\alpha x^2\|_{\mathbb T}
 =\|2\alpha xn+\alpha n^2\|_{\mathbb T}<\varepsilon.
\]
The additional term $2\alpha xn$ depends on $x$ and may cancel
$\alpha n^2$ modulo one. Thus this argument does not force
$\|\alpha n^2\|_{\mathbb T}<\varepsilon$.

\bigskip

With this motivation, we study polynomial configurations whose
shifted-prime parameter satisfies a nil-Bohr restriction. Let
\begin{equation}\label{e:heisenberg-phase}
 Z(n)=\frac{\alpha\beta}{2}n^2-\frac{\alpha\beta}{2}n
       -\alpha n\lfloor\beta n\rfloor.
\end{equation}
We assume that $\alpha$ is badly approximable, meaning that
\begin{equation}\label{e:badly-approximable}
 c_\alpha:=\inf_{m\geq1}m\|m\alpha\|_{\mathbb T}>0.
\end{equation}
We choose a Liouville number $\beta\in(0,1)$ as the limit of an increasing sequence of
rationals $a_j/r_j$ with prime denominators satisfying
\begin{equation}\label{e:horizontal-approximants}
 r_{j+1}\asymp\exp\!\bigl(\exp(\sqrt{\log r_j})\bigr),
 \qquad 0<\beta-\frac{a_j}{r_j}\ll\frac1{r_{j+1}}.
\end{equation}
Given any interval $(u,v)\subset(0,1)$, choose $a_1/r_1$ in its
middle third. Bertrand's postulate and upward rounding give
\eqref{e:horizontal-approximants}, so
$\beta=a_1/r_1+O(r_2^{-1})\in(u,v)$ for sufficiently large $r_1$.
The same bounds give $0<\beta-a_j/r_j=o(r_j^{-D})$ for every
fixed $D>0$. Thus the admissible Liouville numbers are dense
in $(0,1)$.

For an open interval $J\subset\mathbb T$ with $0<|J|<1$, define
\begin{equation}\label{e:explicit-B}
 \mathcal B_{Z,J}:=\{n\in\mathbb Z_{>0}:\{Z(n)\}\in J\},
 \qquad
 \mathcal P_{Z,J}:=\{p\in\mathbb P:p-1\in\mathcal B_{Z,J}\}.
\end{equation}
Write $\pi_{Z,J}(N):=|\mathcal P_{Z,J}\cap[N]|$.
Our aim is to find configurations of the form \eqref{e:intro-pattern} 
inside every sufficiently dense set $A\subset[N]$, with
$p\in\mathcal P_{Z,J}$.

\medskip

Even for the shifts $n,n^2$, an irrational $\beta$ can give
counterexamples to every polylogarithmic density bound. By choosing
successive prime denominators sufficiently large, take
$0<\beta-a_j/r_j<\exp(-r_j^j)$ and set
$N_j=\lfloor\exp(r_j^j/6)\rfloor$. A configuration
$x,x+n,x+n^2$ in $r_j\mathbb Z\cap[N_j]$ has $r_j\mid n$ and
$1\leq n\leq N_j$, so $\lfloor\beta n\rfloor=a_jn/r_j$.
Thus $Z(n)$ is a quadratic polynomial with zero constant term
on these parameters.

Partition $r_j\mathbb Z\cap[N_j]$ using the linear and quadratic
coefficients as torus frequencies, into boxes of side length less
than $1/16$. In a largest cell $A_j$, comparing $x$ with $x+n$
and $x+n^2$ in the respective coordinates gives
$\|Z(n)\|_{\mathbb T}<1/8$. This excludes $J=(1/4,1/2)$, while
$|A_j|/N_j\gg r_j^{-1}$ exceeds $(\log N_j)^{-\kappa}$ for
every fixed $\kappa>0$ and all large $j$.

% Irrationality of $\beta$ alone does not suffice for a polylogarithmic
% density bound. Keep $\alpha$ badly approximable and choose $\beta$
% with prime-denominator approximants satisfying
% $0<\beta-a_j/r_j<\exp(-r_j^j)$, by taking successive prime denominators
% sufficiently large \cite[Book~IX, Proposition~20]{Euclid} and rounding
% upward. Put $N_j=\lfloor\exp(r_j^j/6)\rfloor$ and fix $J=(1/4,1/2)$.
% For $n\leq N_j$ divisible by $r_j$, the identity
% $\lfloor\beta n\rfloor=a_jn/r_j$ gives
% \[
%  Z(n)=\left(\frac{\alpha\beta}{2}-\frac{\alpha a_j}{r_j}\right)n^2
%       -\frac{\alpha\beta}{2}n.
% \]
% Multiply $x$ by the two coefficients in this expression to obtain
% two torus coordinates. Let $A_j$ be a largest cell of
% $r_j\mathbb Z\cap[N_j]$ under the resulting partition into boxes of
% side length less than $1/16$. If $x,x+n,x+n^2\in A_j$ with $n>0$,
% then $r_j\mid n$, and comparing $x$ with $x+n$ in the linear
% coefficient coordinate and with $x+n^2$ in the quadratic coefficient
% coordinate gives $\|Z(n)\|_{\mathbb T}<1/8$. This excludes
% $\{Z(n)\}\in J$, although $|A_j|/N_j\gg 1/r_j$ exceeds
% $(\log N_j)^{-\kappa}$ for every fixed $\kappa>0$ and all large $j$.

\bigskip

To place the restriction \eqref{e:explicit-B} in its natural setting, recall that a basic
polynomial nilsequence has the form
\[
 \psi(n)=F(g(n)\Gamma),
\]
where $G/\Gamma$ is a compact nilmanifold, $G$ is equipped with a
rational filtration, $g:\mathbb Z\to G$ is polynomial for that
filtration, and $F$ is continuous. A set of the form
$\{n:\psi(n)\in U\}$, with $U\subset\mathbb C$ open, is called a
nil-Bohr set; see \cite[Definitions~1.1 and~2.2]{ShaoTeravainen}.
Linear phases give ordinary Bohr conditions, and polynomial phases
give examples on a filtered torus.

The set \eqref{e:explicit-B} comes from the Heisenberg nilmanifold
$H/\Gamma$, where $H=\mathbb R^3$, $\Gamma=\mathbb Z^3$ and the
group law is
\[
 (x,y,z)(x',y',z')=(x+x',\,y+y',\,z+z'+xy').
\]
Left translation by $g_0=(\alpha,\beta,0)$ on $H/\Gamma$ gives
the orbit $g_0^n\Gamma$ of the identity coset, where
$g_0^n=(n\alpha,n\beta,\alpha\beta n(n-1)/2)$. Reducing the second
coordinate by right multiplication subtracts
$n\alpha\lfloor\beta n\rfloor$ from the central coordinate, giving
$\{Z(n)\}$ as the vertical coordinate. Thus $\mathcal B_{Z,J}$
is the return set to the window $0<y<1$, $z\bmod 1\in J$,
which has Haar measure $|J|$.


% In the right quotient, the $n$th iterate of translation by
% $(\alpha,\beta,0)$ from the identity has second horizontal coordinate
% $\{\beta n\}$ and vertical coordinate $\{Z(n)\}$ in the chart
% $0<\{\beta n\}<1$. Thus $\mathcal B_{Z,J}$ is the return set to the
% open window in this chart with vertical coordinate in $J$.
% The window has Haar measure $|J|$.

\bigskip

Peluse \cite{Peluse} obtained quantitative bounds for
distinct-degree polynomial configurations with integer
parameters. For the nonlinear Roth configuration $x,x+n,x+n^2$,
Peluse and Prendiville \cite{PP} proved a polylogarithmic
density bound. Shao and Wang \cite{ShaoWang} obtained
polylogarithmic bounds and a quantitative popular-difference
theorem for general distinct-degree families.
The shifted-prime theorem of \cite{KMTT} gives polylogarithmic
bounds when $n\in\mathbb P-1$. Our first main result retains
that scale while requiring $n\in\mathcal B_{Z,J}$.

% A major advance in quantitative polynomial recurrence was made by
% Peluse \cite{Peluse}, who obtained bounds for distinct-degree
% configurations with integer parameters. Peluse and Prendiville
% \cite{PP} gave a polylogarithmic bound for the nonlinear Roth problem.
% Shao and Wang \cite{ShaoWang}
% subsequently obtained polylogarithmic density bounds, together with
% a quantitative popular-difference theorem. The shifted-prime theorem
% of \cite{KMTT} gives the corresponding scale of density bounds when
% $n\in\mathbb P-1$. Our first main result retains that scale while
% requiring $n\in\mathcal B_{Z,J}$.

\begin{theorem}\label{t:recurrence}
 Assume \eqref{e:fixed-polynomials}, and let $Z$ and $J$ be as in
 \eqref{e:heisenberg-phase}--\eqref{e:explicit-B}. There exists
 $\kappa>0$, depending only on these fixed data, such that, for every
 sufficiently large $N$, every $A\subset[N]$ with
 \[
  |A|\geq N(\log N)^{-\kappa}
 \]
 contains pairwise distinct integers
 \[
  x,\quad x+P_1(p-1),\quad\ldots,\quad x+P_\ell(p-1)
 \]
 for a prime $p\in\mathcal P_{Z,J}$.
\end{theorem}
Writing $n=p-1$, the bracket condition
$\{\alpha p\lfloor\beta p\rfloor\}\in J$ becomes
$\{Z_{\mathrm{br}}(n)\}\in J$, where
\begin{equation}\label{e:bracket-phase}
 Z_{\mathrm{br}}(n)
 :=\alpha(n+1)\lfloor\beta(n+1)\rfloor
 =-Z(n+1)+\frac{\alpha\beta}{2}(n^2+n).
\end{equation}
For $V\in\{Z,Z_{\mathrm{br}}\}$, write
$\mathcal P_{V,J}=\{p\in\mathbb P:\{V(p-1)\}\in J\}$ and
$\pi_{V,J}(N)=|\mathcal P_{V,J}\cap[N]|$.

\begin{corollary}\label{c:bracket-recurrence}
In Theorem~\ref{t:recurrence}, we may replace
$\mathcal P_{Z,J}$ with $\mathcal P_{Z_{\mathrm{br}},J}$.
\end{corollary}

In \eqref{e:bracket-phase}, the shift changes the starting point of
the Heisenberg orbit. At each nonzero Fourier frequency, the
quadratic term is absorbed into the polynomial test, and
multiplication by its exponential preserves the $U^S$ norm for
$S\geq3$.

\bigskip

The recurrence question has an ergodic counterpart. Let $(X,\nu,T)$
be an invertible probability measure-preserving system and let
$\mathbf f=(f_1,\ldots,f_\ell)$, with $f_i\in L^\infty(X)$. Consider
\begin{equation}\label{e:prime-ergodic-average}
 A_N^T(\mathbf f)(x)
 :=\frac1{\pi(N)}\sum_{p\leq N}
       \prod_{i=1}^{\ell}f_i(T^{P_i(p-1)}x).
\end{equation}
The ergodic approach to arithmetic configurations originates in
the work of Furstenberg \cite{Furstenberg} and its polynomial
extension by Bergelson and Leibman \cite{BergelsonLeibman}.
For polynomial averages along the primes, Wooley and Ziegler
\cite{WooleyZiegler} proved $L^2$ convergence. Pointwise convergence
was established by the authors with Krause and  Tao
\cite{KMTTPrimes} for the nonlinear bilinear family $(n,P(n))$,
$\deg P\geq2$, and by Mirek, Wan and Wright
\cite[Theorem~1.1(ii)]{MirekWanWright} for general distinct-degree
families and bounded functions. 
% The latter theorem gives convergence
% of \eqref{e:prime-ergodic-average} by applying it to the polynomials
% $P_i(n-1)$.

Fixed congruence restrictions and linear Bohr conditions can be
incorporated into this pointwise theorem through rotation extensions.
The common step is to establish convergence of the twisted averages
\[
 B_N^\omega(x)
 :=\frac1{\pi(N)}\sum_{p\leq N}e(\omega(p-1))
       \prod_{i=1}^{\ell}f_i(T^{P_i(p-1)}x),
 \qquad e(t):=e^{2\pi it},
\]
for each fixed $\omega\in\mathbb R$. Adjoining the linear iterate $n$
with the constant function $1$ if necessary and relabelling the
polynomials and their inputs, we may assume that $P_1(n)=an$,
with $a\ne0$. On $X\times\mathbb T$, set
$S(x,t)=(Tx,t+\omega/a)$, $G_1(x,t)=f_1(x)e(t)$ and
$G_i(x,t)=f_i(x)$ for $i\geq2$. The corresponding polynomial average
on the product system is exactly $e(t)B_N^\omega(x)$.
The pointwise theorem on this system, followed by Fubini's theorem,
gives almost-everywhere convergence of $B_N^\omega$ for each
fixed $\omega$, including the constant mode
$B_N^0=A_N^T(\mathbf f)$.

Finite Fourier expansion now handles fixed congruence restrictions.
For reduced residue classes, the prime number theorem in arithmetic
progressions also allows normalization by the number of selected
primes. Uniform Fourier approximation similarly gives convergence
for continuous Bohr weights. Continuous approximation and
equidistribution along the primes extend this to interval cutoffs
at irrational frequencies; rational-frequency cutoffs are periodic.
The argument also covers finitely many simultaneous Bohr conditions,
with normalization by their prime count whenever the relative
density is positive.

The restricted limit, however, need not equal the unrestricted one.
For example, consider the \emph{irrational} rotation $Tt=t+\theta$ on
$\mathbb T$, with $f(t)=e(t)$ and $P_1(n)=n$. The unrestricted prime
average tends to zero. For $0<\varepsilon<1/2$, let
$\pi_{\theta,\varepsilon}(N)$ count the primes $p\leq N$ satisfying
$\|\theta(p-1)\|_{\mathbb T}<\varepsilon$. Equidistribution gives
\[
 \frac1{\pi_{\theta,\varepsilon}(N)}
 \sum_{\substack{p\leq N\\
       \|\theta(p-1)\|_{\mathbb T}<\varepsilon}}
       f(T^{p-1}t)
 \longrightarrow
 e(t)\frac{\sin(2\pi\varepsilon)}{2\pi\varepsilon}\ne0.
\]
Thus a linear Bohr restriction can correlate with an eigenfunction
and change the limiting average. The equidistribution input is
classical; see also \cite{ShaoTeravainen}.

Our next theorem shows that the nil-Bohr restriction defining
$\mathcal P_{Z,J}$ preserves the unrestricted limit. The comparison
uses cancellation of the balanced selector against ordinary polynomial
phases of degree at most $d$.

\begin{theorem}\label{t:pointwise}
 Assume \eqref{e:fixed-polynomials}, and let $Z$ and $J$ be as in
 \eqref{e:heisenberg-phase}--\eqref{e:explicit-B}.
 Let $(X,\nu,T)$ be an invertible probability measure-preserving
 system, and let $f_1,\ldots,f_\ell\in L^\infty(X)$. Then
 $\pi_{Z,J}(N)\sim |J|\pi(N)$, and the averages
 \[
  \frac1{\pi_{Z,J}(N)}
   \sum_{\substack{p\leq N\\p\in\mathcal P_{Z,J}}}
        \prod_{i=1}^{\ell}f_i(T^{P_i(p-1)}x)
 \]
 converge almost everywhere to the same limit as
 $A_N^T(\mathbf f)(x)$.
\end{theorem}
The bracket condition from Corollary~\ref{c:bracket-recurrence}
also preserves the limit of the unrestricted prime averages.
\begin{corollary}\label{c:bracket-pointwise}
In Theorem~\ref{t:pointwise}, we may replace $\mathcal P_{Z,J}$
and $\pi_{Z,J}(N)$ with $\mathcal P_{Z_{\mathrm{br}},J}$
and $\pi_{Z_{\mathrm{br}},J}(N)$, respectively.
\end{corollary}





\bigskip

A broader question is which nilsequence weights can be inserted into
a discrete configuration count or a polynomial ergodic average while
preserving its main term. The Heisenberg selector gives a concrete
case in which the full Gowers norms detect structure, yet the
correlations with ordinary polynomial phases are small on the
progressions used in the density increment. An extension to general
nilsequences would require a quantitative hypothesis on these
correlations that remains effective under the arithmetic rescalings. We do not pursue this generalization in this paper. 

One can also ask whether, for every $k\geq2$, there is a progression
$p,p+r,\ldots,p+(k-1)r$, with $r>0$, entirely contained in
$\mathcal P_{Z,J}$. For fixed $J$, this follows directly from
$\pi_{Z,J}(N)\sim |J|\pi(N)$ and the relative-density form of the
Green--Tao theorem \cite{GreenTao}. If $J=J_N$ shrinks with $N$, then for $k\geq4$ the quantitative
Green--Tao theorem of the second author and Wang
\cite[Theorem~1.1]{TeravainenWang} gives a $k$-term progression
among the selected primes up to $N$ for all sufficiently large
$N$ whenever
\[
 \frac{\pi_{Z,J_N}(N)}{\pi(N)}
 > C_k\exp\bigl(-(\log\log\log N)^{c_k}\bigr)
\]
for suitable constants $C_k,c_k>0$ depending only on $k$.

On the ergodic side, one can ask how far the pointwise comparison
extends to inputs in $L^p$. Where a suitable maximal estimate is available for the unrestricted
prime averages, domination gives the corresponding estimate for
the restricted averages. The endpoint $p=1$ already poses
difficulties for the unrestricted averages.
A question specific to the comparison is whether one can bound the
variation of the difference between the restricted and unrestricted
averages. We do not pursue this direction in this paper either. 

The arithmetic argument also suggests studying other weights. A
unified result for a class of multiplicative weights would require
conditions on their multiplicative structure that yield cancellation
against vertical modes of the selector and ordinary polynomial
phases on the progressions used in the proof. 

\subsection{Overview of the proofs}

For simplicity, consider the configuration $x,x+n,x+n^2$, with
$n=p-1$ and $\{Z(n)\}\in J$. For $A\subset[N]$ with $|A|/N =\delta$ and
$M=\lfloor N^{1/2}\rfloor$, let
\begin{equation}\label{e:formprim}
 \mathcal C(F)=\frac1{NM}\sum_{x\in\mathbb Z}\sum_{1\leq n\leq M}
 \Lambda(n+1)F(Z(n))1_A(x)1_A(x+n)1_A(x+n^2).
\end{equation}
We seek $\mathcal C(1_J)=|J|\mathcal C(1)+o(1)$ uniformly in $A$.
The prime number theorem gives $\mathcal C(1)\ll1$, so smooth
majorants and minorants of $1_J$ with means approaching $|J|$
reduce the comparison to $\mathcal C(F-\mu_F)=o(1)$ for smooth
$F$, where $\mu_F=\int_{\mathbb T}F$. Fourier approximation then
leaves the counts $\mathcal C(e(h\,\cdot))$ with $h\ne0$.

The density increment needs a comparison after each rescaling.
Fix a nonnegative smooth minorant $F$ of $1_J$ with $\mu_F>0$.
Restricting $x$ to a progression of step $q$ and taking $n=qy$
gives the shifts $y,qy^2$ at spatial scale $qM^2$, where $M$ is
now the length of the $y$-range. The weight becomes $\Lambda_qF_q$,
where $\Lambda_q(y)=\varphi(q)\Lambda(qy+1)/q$ and
$F_q(y)=F(Z(qy))$. At the smooth moduli in the iteration, the
prime model for $\Lambda_q$ is $1$. To compare the selected-prime
count with $\mu_F$ times the unweighted count, write
\[
 \Lambda_qF_q-\mu_F=(\Lambda_q-1)F_q+(F_q-\mu_F).
\]
At density $\delta$, the unweighted count for the constant-density
inputs is of order $\delta^3$. We need an error small compared
with $\delta^3$, uniformly over the moduli in the iteration.
The two natural candidates for controlling these weighted counts
are generalized von Neumann estimates using full Gowers norms
and estimates using correlations with ordinary polynomial phases.

Full Gowers norms give such an estimate
\cite[Theorem~4.2]{Teravainen}, and the version in
\cite[Proposition~3.3]{KMTT} accommodates the rescaled spatial
scale. We use it for $(\Lambda_q-1)F_q$. To insert a Fourier mode
of $F_q$, we localize it to a short interval in the horizontal
coordinate $\{\beta qy\}$. On each cube where the cutoff is nonzero
at every vertex, the vertical phase is quadratic in the cube
variables and cancels in the $U^S$ product for $S\geq3$. The
high-order Gowers bound for $\Lambda_q-1$ thus controls the count
with the selector inserted.

For the centered selector $F_q-\mu_F$, the same cube calculation gives a
lower bound. Already at $q=1$, equidistribution gives a nonzero
correlation with a localized vertical mode. More precisely, for
$h\ne0$ with $\widehat F(h)\ne0$ and a nonzero nonnegative smooth
cutoff $\chi$ supported on a short interval inside $(0,1)$,
\[
 0<|\widehat F(h)|\int_{\mathbb T}\chi
 \ll \|\widehat\chi\|_{\ell^1}
       \|F(Z(\cdot))-\mu_F\|_{U^3[M]}
\]
for all sufficiently large $M$. By monotonicity, full Gowers
norms give no decay for the centered selector at any order
$S\geq3$. A general weighted bound using $U^2$ also fails.
For irrational $\theta$, the weight $e(-\theta n^2)$ has small
$U^2[M]$, yet the spatial phases $e(-\theta x)$ and
$e(\theta(x+n^2))$ cancel it exactly on the supporting interval.

Second author's fixed-family estimate
\cite[Theorem~4.1]{Teravainen} already uses local linear phases to
obtain polynomial-phase control. Appendix~\ref{s:why-Cq}
constructs growing smooth moduli at which the $U^3$ obstruction
persists and lifting (dilating) the count back to the fixed shifts $n,n^2$
gives no decay through estimate \cite[Theorem~4.1]{Teravainen}. This motivates
our local generalized von Neumann theorem, which tests the
weight on parameter progressions and applies directly to
$y,qy^2$ with constants uniform in $q$. After differencing the
spatial inputs, we use the coefficient-uniform inverse estimate
of \cite[Lemma~3.4]{KMTT}, based on Shao and Wang
\cite[Theorem~2.4]{ShaoWang}, to extract linear phases on spatial
progressions. Evaluating these phases along the shifts gives a
quadratic phase on each parameter progression. For a general
distinct-degree family, the resulting phase has degree at most
the largest degree $d$.

For the ambient count $\mathcal C(F-\mu_F)$, the same cube argument
controls the contribution of $\Lambda-\Lambda_{\Cramer,\widetilde Q}$
to the quadratic-phase tests. The Cram\'er cutoff $\widetilde Q$
is justified by the explicit zero-free bound in \cite{BanksSiegel}.
Expanding the sieve weights in the Cram\'er term gives tests on
progressions. Divisors containing the approximating prime $r_j$
have negligible total contribution. The remaining estimate, both
for the Cram\'er term and for the centered selector at the smooth
moduli, is
\[
 \sup_{\deg\phi\leq2}
 \left|\mathbb E_{n\in P}
     \bigl(F(Z(n))-\mu_F\bigr)e(-\phi(n))\right|
 \ll_{F,D}(\log M)^{-D}
 \qquad(D>0).
\]
Here $P$ runs over the long progressions from the sieve or,
after writing $n=qy$, from the local generalized von Neumann
estimate. Their steps are coprime to $r_j$.

To prove this estimate, we adjoin a circle carrying the test
phase and apply Green and Tao's quantitative equidistribution
theorem \cite{GreenTaoNil}. A large correlation forces an integer
multiple of the test phase to agree with a horizontal linear
phase up to an integer-valued polynomial and a slowly varying
error. After dividing by that integer, the polynomial factor
is constant on residue classes and the error is handled by
partial summation. The horizontal twist is absorbed on a finite
cover. The localized vertical mode still has zero mean on each
vertical fibre there, and the horizontal Diophantine bounds give
the required cancellation.

At the ambient scale, the resulting prime-weighted phase bound
and Second author's estimate give
$$\mathcal C(F-\mu_F)=\mathcal C(F)-\mu_F\mathcal C(1)=o(1).$$
At the smooth moduli, the local theorem gives the required
saving for the count with weight $F_q-\mu_F$, completing the
comparison with the unweighted count. If $A$ contains no
selected-prime configuration, this comparison makes its
unweighted count smaller than that for the constant-density
inputs by a fixed multiple of $\delta^3$. The weak regularity
argument of \cite{KMTT} then gives a long progression on which
the density increases by a fixed factor. A selected-prime
configuration in the rescaled set would lift to one in $A$,
so the uniform comparison lets us repeat the increment until
the density exceeds one, or $\delta=o(\log N)^{-\kappa}$.

For pointwise convergence the polynomial family stays fixed.
The prime-weighted phase bound at $q=1$ and \cite[Theorem~4.1]{Teravainen} give an $L^1$ comparison summable along an increasing subsequence whose consecutive
ratios tend to one.
A gap estimate extends the comparison to all scales. Smoothing
passes to the interval, and taking all inputs equal to one
gives the normalization by the number of selected primes.
The unrestricted prime averages \cite{MirekWanWright} then identify the common limit.

\begin{remark}
At the moduli $q$ used in the density increment, which are divisible
by $\operatorname{lcm}(1,\ldots,w)$, the bound
$|T_{N,M,F_q-\mu_F}(\mathbf f)|\ll w^{-\eta}$ needed for the
density increment can be proved directly using the fixed local
factors from the proof of \cite[Lemma~5.3]{KMTT}. The argument
uses the same spatial differencing, duality, and extraction of
local linear phases as the proof of Theorem~\ref{t:main}.
Here the divisibility of $q$ absorbs the auxiliary denominators
and eliminates the pigeonholing over major-arc frequencies.
Theorem~\ref{t:main} gives polynomial-phase control for arbitrary
bounded weights over a wider range of moduli.
\end{remark}

\subsection{Organization}

Section~\ref{s:recurrence} reduces Theorem~\ref{t:recurrence}
to the density increment in Proposition~\ref{p:selected-density-increment}
and derives that increment from the smooth-modulus comparison
in Theorem~\ref{t:relative-selector-comparison}.
Section~\ref{s:relative} proves the comparison using the local
generalized von Neumann estimate in Theorem~\ref{t:main} and
the phase estimates in Proposition~\ref{p:heisenberg-phase}.
Their proofs are given in Sections~\ref{s:local-vn}
and~\ref{s:example}, respectively. Section~\ref{s:example}
also proves the prime-weighted phase estimate used at $q=1$
for the pointwise comparison in Section~\ref{s:ergodic}.
Appendix~\ref{s:why-Cq} compares the local estimate with the
bound obtained by dilation to fixed shifts.

\section{Preliminaries}\label{s:preliminaries}

% Each density increment restricts the spatial variable to a progression.
% We fix notation for the resulting polynomial coefficients, counting
% forms, and weights, so that the argument can be repeated after
% rescaling. Two norms enter the counting comparison. Full Gowers
% norms control the discrepancy of the prime weight from its model,
% whereas correlations with polynomial phases on progressions control
% the centered selector. We also record the nilmanifold notation used
% to estimate those correlations.

Write \([N]=\{1,\ldots,N\}\), \(\e(t)=e^{2\pi i t}\), and
\(\E_{n\in P}=|P|^{-1}\sum_{n\in P}\) for a finite nonempty set \(P\).
Unrestricted sums over \(x\) are over \(\Z\).  A function is
\(1\)-bounded if its modulus is at most \(1\).
We write $\pi(N)$ for the number of primes up to $N$ and
$\varphi(q)$ for Euler's totient function. The von Mangoldt
function is $\Lambda(p^j)=\log p$ for primes $p$ and $j\geq1$,
and is zero otherwise. Our Fourier conventions are
\[
 \widehat f(\xi)=\sum_{x\in\mathbb Z}f(x)\e(-x\xi),
 \qquad
 \widehat F(h)=\int_{\mathbb T}F(t)\e(-ht)\,dt.
\]

A filtration of degree at most \(s\) on a connected, simply connected
nilpotent Lie group \(G\) is a descending sequence
\[
 G=G_0=G_1\supseteq G_2\supseteq\cdots\supseteq G_{s+1}=\{e\},
 \qquad [G_i,G_j]\subseteq G_{i+j}.
\]
If \(\Gamma\le G\) is discrete and cocompact and the filtration is
rational with respect to \(\Gamma\), we call \(G/\Gamma\), equipped with
this filtration, a filtered nilmanifold.  A sequence \(g:\Z\to G\) is
polynomial when
\[
 \partial_{h_1}\cdots\partial_{h_i}g(n)\in G_i
 \qquad(i\ge1),
\]
where \(\partial_hg(n)=g(n+h)g(n)^{-1}\).  A basic nilsequence is a
sequence \(F(g(n)x_0)\), where \(x_0\in G/\Gamma\) and
\(F:G/\Gamma\to\C\) is continuous.  We use the standard quantitative
notion of complexity arising from a rational Mal'cev basis
\cite[Section~2]{GreenTaoNil}.  All
nilmanifolds occurring in the applications are fixed, so their complexity
is absorbed in the implicit constants.

For a finitely supported function \(g:\mathbb Z\to\mathbb C\), use the
unnormalized Gowers norm
\[
 \|g\|_{U^s(\mathbb Z)}^{2^s}
 :=\sum_{x,h_1,\ldots,h_s\in\mathbb Z}
   \prod_{\omega\in\{0,1\}^s}
   \mathcal C^{|\omega|}g(x+\omega\cdot h),
\]
where \(\mathcal C\) denotes complex conjugation.  On a finite nonempty interval or arithmetic progression
\(I\), functions are extended by zero and we normalize as in
\cite[Section~2.3]{KMTT}:
\begin{equation}\label{e:normalized-gowers}
 \|g\|_{U^s[I]}
 :=\frac{\|g\mathbf1_I\|_{U^s(\mathbb Z)}}
          {\|\mathbf1_I\|_{U^s(\mathbb Z)}}.
\end{equation}
In particular, \(\|g\|_{U^1[I]}=|\mathbb E_{n\in I}g(n)|\), and
normalized Gowers monotonicity gives
\(\|g\|_{U^r[I]}\ll_{r,s}\|g\|_{U^s[I]}\) when \(1\le r\le s\).

For a finite nonempty arithmetic progression \(P\), define
\begin{equation}\label{e:u-def}
 \|w\|_{u^s[P]}:=\sup_{\phi\in\mathbb R[y],\ \deg\phi<s}
 \left|\mathbb E_{n\in P}w(n)\e(-\phi(n))\right|.
\end{equation}
Repeated Cauchy--Schwarz gives
\begin{equation}\label{e:u-below-U}
 \|w\|_{u^s[P]}\ll_s\|w\|_{U^s[P]}.
\end{equation}
Following~\cite[Definition~2.3]{ShaoWang}, we use the coefficient class
produced by arithmetic rescaling.

\begin{definition}[$(C,q)$-coefficients]\label{d:Cq}
Let $C\ge1$ and $q\in\mathbb N$.  A polynomial
\[
 P(y)=a_ey^e+\cdots+a_1y\in\Z[y],\qquad e\ge1,
\]
with $P(0)=0$ has \emph{$(C,q)$-coefficients} if
\[
 |a_j|\le C|a_e|\quad(1\le j<e),
 \qquad a_e=a'_eq^{e-1}
\]
for some $a'_e\in\Z$ with $0<|a'_e|\le C$.
\end{definition}
% The leading coefficient is therefore allowed to carry the power $q^{e-1}$ produced by the rescaling $P(y)\mapsto P(qy)/q$.  This dependence on $q$ is not covered uniformly by a bounded-coefficient theorem.

For a family $\mathbf R=(R_1,\ldots,R_\ell)$ with $R_0=0$ and
$\mathbf f=(f_0,\ldots,f_\ell)$, write
\begin{equation}\label{e:Tdef}
 T_{N,M,v}^{\mathbf R}(\mathbf f)
 :=\frac1{NM}\sum_{x\in\mathbb Z}\sum_{1\le y\le M}
 v(y)\prod_{i=0}^\ell f_i(x+R_i(y)).
\end{equation}
We omit the superscript when the family is clear. In a count with
$k$ input functions we write $k=\ell+1$. We use
\[
 \Lambda_q(y)=\frac{\varphi(q)}q\Lambda(qy+1),\qquad
 F_q(y)=F(Z(qy)),\qquad \mu_F=\int_{\mathbb T}F.
\]
For $w\ge2$ an integer, let $\mathcal L(w)=\operatorname{lcm}(1,\ldots,w)$.
Let $P^+(q)$ denote the largest prime factor of $q$, with $P^+(1)=1$.

% In the rescaled count, we restrict the spatial variable to
% $x=a+qu$ and the parameter to $n=qy$. Then
% $x+P_i(n)=a+q(u+P_i(qy)/q)$, so the shifts in the new spatial
% variable are $P_{i,q}(y)=P_i(qy)/q$. They have
% $(C,q)$-coefficients for a fixed $C$ depending on the original
% family. If the next progression has step $q^B$ in the rescaled
% variable, the new modulus is $q'=q^{B+1}$ and
% $q^B P_{i,q'}(y)=P_{i,q}(q^By)$. The prime is preserved under
% this change of variables, since $q'y+1=q(q^By)+1$.



\section{Reduction of recurrence to a counting comparison}\label{s:recurrence}

Suppose a dense set contains no configuration of
Theorem~\ref{t:recurrence}. We seek a long progression on which the
set has higher density and on which the same absence of configurations
remains after rescaling. We state the density increment first and
show how it can be iterated. We then derive it by comparing the
selected-prime count with the unweighted count.

Write \(k=\ell+1\) and \(r=\deg P_1\). At a stage with modulus \(q\),
the shifts are \(P_{j,q}(y)=P_j(qy)/q\), and the parameter must satisfy
\(qy+1\in\mathcal P_{Z,J}\). The coefficient of $y^j$ in $P_{i,q}(y)$ is $q^{j-1}$
times that in $P_i$, so the rescaled family has $(C,q)$-coefficients
for a fixed $C$ depending only on the original family.

\begin{proposition}[Density increment]\label{p:selected-density-increment}
There are constants \(c,\eta,K,E,C_1>0\) and \(B_*\in\mathbb N\),
depending only on the fixed data of Theorem~\ref{t:recurrence}, with the
following property for sufficiently large \(N\). Let \(A\subset[N]\)
have density \(0<\delta\le1\), and suppose \(q,w\in\mathbb N\) satisfy
\begin{equation}\label{e:increment-hypotheses}
 K\delta^{-E}\le w\le(\log N)^c,\qquad
 \mathcal L(w)\mid q,\qquad P^+(q)\le w,\qquad \log q\le w^2.
\end{equation}
Put \(M=\lfloor(N/q^{d-1})^{1/d}\rfloor\). Suppose that \(A\) contains
no configuration of pairwise distinct points
\(x,x+P_{1,q}(y),\ldots,x+P_{\ell,q}(y)\) with
\(y\ge1\) and \(qy+1\in\mathcal P_{Z,J}\).
Then there are \(1\le B\le B_*\), \(a\in\mathbb Z\), and \(N'\in\mathbb N\)
such that
\begin{equation}\label{e:increment-output}
 a+q^B[N']\subset[N],\qquad
 |A\cap(a+q^B[N'])|\ge(1+\eta)\delta N',\qquad
 N'\gg q^{-C_1}\delta^{C_1}M^r.
\end{equation}
\end{proposition}

To deduce Theorem~\ref{t:recurrence} from the density increment, we
follow the iteration in \cite[Section~5, proof of Theorem~1.2]{KMTT}.
We first pass to a densest residue class modulo an initial smooth
modulus and rescale. Each application of the proposition gives a
denser progression whose rescaled set still lacks the required
configuration. We track the lengths and moduli to ensure that the
proposition continues to apply until the density would exceed one.

\begin{proof}[Proof of Theorem~\ref{t:recurrence} assuming
Proposition~\ref{p:selected-density-increment}]
Suppose \(A\subset[N]\) contains no required configuration and has
density \(\delta\ge(\log N)^{-\kappa}\). Take the constants in
Proposition~\ref{p:selected-density-increment}, and set
\(w=\lceil K\delta^{-E_1}\rceil\), where \(E_1\ge E\) will be fixed below.
Set \(q_0=\mathcal L(w)\). Among the nonempty residue classes
modulo \(q_0\) in \([N]\), choose \(u_0+q_0[N_0]\) maximizing
\(\lvert A\cap(u_0+q_0[N_0])\rvert/N_0\). Define
\[
 A_0=\{x\in[N_0]:u_0+q_0x\in A\}.
\]
By averaging over the residue classes, \(\delta_0=|A_0|/N_0\ge\delta\).
For sufficiently small \(\kappa>0\) we have
\[\log q_0= \log\mathcal L(w)=\sum_{m\le w}\Lambda(m)\ll w=o(\log N).\] 
Thus the progression
has length \(N_0\ge N/q_0-1\ge N/(2q_0)\).  A configuration
in \(A_0\) with shifts \(P_{j,q_0}(y)\) and \(q_0y+1\in\mathcal P_{Z,J}\)
would lift under the affine change of variables to a required
configuration in \(A\).

Whenever Proposition~\ref{p:selected-density-increment} applies at
\(A_i,N_i,q_i\), it gives a progression \(a_i+q_i^{B_i}[N_{i+1}]\).
Define
\[
 A_{i+1}=\{x\in[N_{i+1}]:a_i+q_i^{B_i}x\in A_i\}.
\]
Then \(\delta_{i+1}=|A_{i+1}|/N_{i+1}\ge(1+\eta)\delta_i\).
Set \(q_{i+1}=q_i^{B_i+1}\). The identity
\[
 q_i^{B_i}P_{j,q_{i+1}}(y)=P_{j,q_i}(q_i^{B_i}y),
 \qquad q_{i+1}y+1=q_i(q_i^{B_i}y)+1
\]
preserves both the shifts and the selected prime under the lift
\(x\mapsto a_i+q_i^{B_i}x\). The hypotheses concerning configurations
therefore persist. Each \(q_i\) is a power of \(q_0\), so
\(\mathcal L(w)\mid q_i\) and \(P^+(q_i)\le w\). The lower bound
\(w\ge K\delta_i^{-E}\) follows from \(\delta_i\ge\delta\).

It remains to check the two upper bounds in
\eqref{e:increment-hypotheses} through
\(i_*=\lfloor\log(\delta^{-1})/\log(1+\eta)\rfloor+1\) increments.
Iterating \(q_{h+1}=q_h^{B_h+1}\), we obtain, for every
\(i\le i_*\) reached by the iteration,
\[
 \log q_i
 =\left(\prod_{h=0}^{i-1}(B_h+1)\right)\log q_0
 \le (B_*+1)^i\log q_0
 \ll \delta^{-O(1)}w.
\]
The exponent here depends only on \(\eta,B_*\). Fix \(E_1\) larger than
this exponent and increase \(K\) by a fixed factor to ensure
\(\log q_i\le w^2\). From \eqref{e:increment-output}, with
\(M_i=\lfloor(N_i/q_i^{d-1})^{1/d}\rfloor\), we have
\[
 \log N_{i+1}\ge \frac rd\log N_i
    -O\bigl(\log q_i+\log(\delta^{-1})+1\bigr).
\]
Since \(\log N_0\ge\log N-O(w)\) and
\((r/d)^i\gg \delta^{O(1)}\) for \(i\le i_*\), iteration gives
\[
 \log N_i\ge c_2\delta^{C_2}\log N
 -O\bigl((1+\log(\delta^{-1}))(w^2+\log(\delta^{-1}))\bigr)
\]
for some fixed \(c_2,C_2>0\).  Taking \(\kappa\) sufficiently small after \(E_1\), it gives
\(\log N_i\ge(\log N)^{3/4}\) and
\(w\le(\log N_i)^c\) at every stage. We may therefore carry out all
\(i_*\) increments. After \(i_*\) increments the density would be at least
\((1+\eta)^{i_*}\delta>1\), which is impossible. Hence some \(A_i\) with \(i<i_*\) contains
a configuration with shifts \(P_{j,q_i}(y)\) and
\(q_i y+1\in\mathcal P_{Z,J}\). The successive affine maps lift this
configuration back to \(A\).
\end{proof}

For the bracket condition, we apply Theorem~\ref{t:recurrence}
within one cell of a finite partition of $A$ that controls
$Z_{\mathrm{br}}(n)+Z(n)$ modulo one.

\begin{proof}[Proof of Corollary~\ref{c:bracket-recurrence}]
Choose an open interval $J_0$ with $-\overline{J_0}\subset J$.
Take $\varepsilon>0$ small enough that
\[
 \varepsilon<\min(\beta,1-\beta),\qquad
 (3+|\alpha|)\varepsilon
 <\operatorname{dist}(-\overline{J_0},\mathbb T\setminus J).
\]
Enlarge the polynomial family to
\[
 \{P_1,\ldots,P_\ell\}
 \cup
 \bigl\{n^j:\ j\in\{1,2\},\
       j\notin\{\deg P_i:1\le i\le\ell\}\bigr\}.
\]
Order this family by degree and write its linear and quadratic
members as $an$ and $bn^2+cn$, respectively, with $ab\ne0$, and put $\theta=\alpha\beta/2$.

Partition $A$ according to the positions of
\[
 \left(
 \frac{\alpha x}{a},\,
 \frac{\beta x}{a},\,
 \frac{\theta(1-c/b)x}{a},\,
 \frac{\theta x}{b}
 \right)\pmod 1
\]
in fixed boxes of side length less than $\varepsilon$. Apply
Theorem~\ref{t:recurrence} with interval $J_0$ to a largest cell
and the enlarged family. The cell has size $\gg |A|$, so this
fixed loss is absorbed by choosing a smaller $\kappa$. The theorem
gives a parameter $n=p-1$ with $\{Z(n)\}\in J_0$. The points $x$, $x+an$, and $x+bn^2+cn$ lie in the same cell,
so each coordinate difference has torus norm less than
$\varepsilon$. Comparing $x+an$ with $x$ in the first three
coordinates gives
\[
 \|\alpha n\|_{\mathbb T}<\varepsilon,\qquad
 \|\beta n\|_{\mathbb T}<\varepsilon,\qquad
 \left\|\theta\left(1-\frac cb\right)n\right\|_{\mathbb T}
 <\varepsilon,
\]
since these coordinates multiply $x$ by $\alpha/a$, $\beta/a$,
and $\theta(1-c/b)/a$, respectively. Comparing $x+bn^2+cn$
with $x$ in the fourth coordinate gives
\[
 \left\|\frac{\theta}{b}(bn^2+cn)\right\|_{\mathbb T}
 =\left\|\theta n^2+\frac{\theta c}{b}n\right\|_{\mathbb T}
 <\varepsilon.
\]
The third and fourth differences add to $\theta(n^2+n)$, so
the triangle inequality gives
\[
 \|\theta(n^2+n)\|_{\mathbb T}
 \leq
 \left\|\theta\left(1-\frac cb\right)n\right\|_{\mathbb T}
 +\left\|\theta n^2+\frac{\theta c}{b}n\right\|_{\mathbb T}
 <2\varepsilon.
\]

Let $m$ be the nearest integer to $\beta n$. Since
$|\beta n-m|<\varepsilon<\min(\beta,1-\beta)$, we have
$\lfloor\beta(n+1)\rfloor=m$. Also
$m-\lfloor\beta n\rfloor\in\{0,1\}$. The definitions give
\[
 Z_{\mathrm{br}}(n)+Z(n)
 =\theta(n^2+n)
   +\alpha n\bigl(m-\lfloor\beta n\rfloor\bigr)
   -\alpha(\beta n-m).
\]
Thus $\|Z_{\mathrm{br}}(n)+Z(n)\|_{\mathbb T}
<(3+|\alpha|)\varepsilon$. Since $\{Z(n)\}\in J_0$, the choice
of $\varepsilon$ gives $\{Z_{\mathrm{br}}(n)\}\in J$.
Keeping the original shifts gives the desired configuration.
\end{proof}

% \begin{proof}[Proof of Corollary~\ref{c:bracket-recurrence}]
% Apply Proposition~\ref{p:selected-density-increment} with
% $V=Z_{\mathrm{br}}$. The comparison used to prove the proposition
% now counts primes $qy+1$ satisfying
% $\{\alpha(qy+1)\lfloor\beta(qy+1)\rfloor\}\in J$.
% In the lift from modulus $q'=q^{B+1}$ to modulus $q$, the parameter
% $y$ becomes $q^By$. The prime $q'y+1=q(q^By)+1$ is preserved,
% so its bracket condition is preserved as well. The iteration in
% the proof of Theorem~\ref{t:recurrence} therefore gives the claimed
% density bound and configuration.
% \end{proof}

We now turn to the counting comparison that gives the density
increment. The increment requires the count with weight
\(\Lambda_qF_q\) to be close to the count with constant weight
\(\mu_F\), uniformly over the smooth moduli reached in the iteration.

Let \(c_0\in(0,1/11)\) be the fixed constant in
\cite[Definition~2.3]{KMTT}, and put
\(\widetilde Q_N=\lceil\exp((\log N)^{c_0})\rceil\).
For \(w_0\) and a sufficiently small \(c_{\rm rel}>0\) supplied by
the theorem below, we impose
\begin{equation}\label{e:relative-smooth-hypotheses}
 w_0\le w\le(\log\widetilde Q_N)^{c_{\rm rel}},\qquad
 8\prod_{p\le w}p\mid q,\qquad P^+(q)\le w,\qquad \log q\le w^2.
\end{equation}
In particular, \(\log q\le(\log\widetilde Q_N)^{2c_{\rm rel}}\), so
the choice of \(c_{\rm rel}\) keeps \(q\) within the range of the
prime-model estimates.
We use the counting form $T$ defined in \eqref{e:Tdef}.

\begin{theorem}\label{t:relative-selector-comparison}
Fix $C,d,k$, the parameters of $Z$, and a Lipschitz function
$F:\mathbb T\to\mathbb C$. There are $c_{\rm rel},\eta_{\rm sel}>0$
and an integer $w_0\geq16$ such that the following holds for all
sufficiently large integers $N$. Let $q,w\in\mathbb N$ satisfy
\eqref{e:relative-smooth-hypotheses}, and put
$M=\lfloor(N/q^{d-1})^{1/d}\rfloor$. Let $P_1,\ldots,P_{k-1}$
have distinct increasing degrees, largest degree $d$, and
$(C,q)$-coefficients. Then for all $1$-bounded
$f_0,\ldots,f_{k-1}$ supported on $[N]$,
\begin{equation}\label{e:relative-selector-constant}
 T_{N,M,\Lambda_q F_q}(\mathbf f)
 =\mu_F T_{N,M,1}(\mathbf f)+O(w^{-\eta_{\rm sel}}).
\end{equation}
The implied constant is uniform over the stated coefficient class.
\end{theorem}

Theorem~\ref{t:relative-selector-comparison} is proved in
Section~\ref{s:relative}. Fix a Lipschitz function \(0\le F\le 1\)
supported in \(J\), with \(\mu_F>0\), and write \(T_v=T_{N,M,v}\)
for the family with shifts \(P_{j,q}\). If \(A\) contains no
required configuration, the comparison places its unweighted
count below that of the constant-density model
\(\delta\mathbf 1_{[N]}\). Weak regularity then gives a long
progression on which \(A\) has density at least \((1+\eta)\delta\).

\begin{proof}[Proof of Proposition~\ref{p:selected-density-increment}
assuming Theorem~\ref{t:relative-selector-comparison}]
Fix \(C\) large enough that all the \(P_{j,q}\) have
\((C,q)\)-coefficients. Take \(c>0\) sufficiently small in \eqref{e:increment-hypotheses}.
Then 
\[\log q\le(\log N)^{2c}=o(\sqrt{\log N}),\] so
\(M=N^{1/d-o(1)}\).
Since \(w\ge K\ge16\), the divisibility by \(\mathcal L(w)\) implies
\(8\prod_{p\le w}p\mid q\). With \(c\) small enough and \(K\) large
enough, \eqref{e:relative-smooth-hypotheses} holds as well.
Thus the comparison \eqref{e:relative-selector-constant} applies.

For $1\leq y\leq\tau M$ and $e=\deg P_j$, the coefficient bounds
give $|P_{j,q}(y)|\ll_{C,d}q^{e-1}(\tau M)^e\ll_{C,d}\tau N$.
Choose $\tau>0$ small enough that all shifts are at most $N/20$
in absolute value. For each such $y$, at least $N/2$ values of
$x$ have all shifted points in $[N]$. Hence, for some fixed $c_3>0$,
\[T_1(\delta\mathbf1_{[N]},\ldots,\delta\mathbf1_{[N]})
\ge c_3\delta^k.\] 
For \(\mathbf f=(\mathbf1_A,\ldots,\mathbf1_A)\), the only possible
contributions to \(T_{\Lambda_qF_q}(\mathbf f)\) come from higher prime
powers and from coincident shifts. Including $P_{0,q}=0$, each difference $P_{i,q}-P_{j,q}$ with
$i\ne j$ is nonzero, so there are $O_{d,k}(1)$ values of $y$
with coincident shifts. The total contribution of both types is
\[
 \ll_{d,k}\frac{\log(qM+2)}{M}
       +\frac{(qM+1)^{1/2}(\log(qM+2))^2}{M}
 \ll N^{-1/(3d)}
\]
for large \(N\). Here we used \(0\le F\le1\), the bound
\(\Lambda_q(y)\le\Lambda(qy+1)\), and the elementary count of higher
prime powers up to \(qM+1\). Since \(\delta^{-E}\ll(\log N)^c\)
by \eqref{e:increment-hypotheses}, this contribution is \(o(\delta^k)\).
Choose \(E\) with \(E\eta_{\rm sel}\ge k\) and \(K\) large enough that
the comparison error is at most \(\mu_F c_3\delta^k/4\).
Dividing the comparison for \(\mathbf f\) by \(\mu_F\) now gives
\begin{equation}\label{e:unweighted-discrepancy-hyp}
 T_1(\delta\mathbf1_{[N]},\ldots,\delta\mathbf1_{[N]})
 -T_1(\mathbf1_A,\ldots,\mathbf1_A)\ge \tfrac12 c_3\delta^k.
\end{equation}

Apply the weak regularity lemma \cite[Lemma~5.3]{KMTT} to the
unweighted form with accuracy \(\varepsilon=c_3\delta^k/8\).
Increasing \(E,K\) ensures that the required
\(\mathcal L(\lceil\varepsilon^{-A'}\rceil)\mid q\) follows from
\(\mathcal L(w)\mid q\). Its small-scale alternative is impossible:
\(M=N^{1/d-o(1)}\) whereas \((q/\varepsilon)^{O_{C,d,k}(1)}=N^{o(1)}\).
Thus we obtain local factors of \([N]\) with common modulus \(q^B\),
\(1\le B\le B_*\), whose atoms have lengths
\(\gg q^{-C_1}\delta^{C_1}M^r\) for a suitable fixed \(C_1\).
Let \(H_j\) be the conditional expectation of \(\mathbf 1_A\) on
the \(j\)th factor, so that \(H_j(x)=|A\cap I|/|I|\) on each
atom \(I\) of that factor. Extend each $H_j$ by zero off $[N]$. Weak regularity gives
\[
 \bigl|T_1(H_0,\ldots,H_\ell)
       -T_1(\mathbf 1_A,\ldots,\mathbf 1_A)\bigr|
 \le \varepsilon = c_3\delta^k/8.
\]
Combining this with \eqref{e:unweighted-discrepancy-hyp}, we obtain
\begin{align}\label{e:diff1a}
    T_1(\delta\mathbf 1_{[N]},\ldots,\delta\mathbf 1_{[N]})
 -T_1(H_0,\ldots,H_\ell)
 \ge \tfrac12 c_3\delta^k-\varepsilon
 = \tfrac38 c_3\delta^k.
\end{align}
Suppose $H_j\leq(1+\eta)\delta$ on $[N]$ for every $j$.
Since $H_j$ and $\delta\mathbf1_{[N]}$ both have mass $\delta N$,
\[
 \|H_j-\delta\mathbf1_{[N]}\|_{\ell^1}
 =2\sum_{x\in[N]}(H_j(x)-\delta)_+
 \leq 2\eta\delta N.
\]
Telescope by replacing one input at a time. The $j$th summand
satisfies
\[
\begin{aligned}
 &\bigl|T_1(H_0,\ldots,H_{j-1},H_j-\delta\mathbf 1_{[N]},\delta\mathbf 1_{[N]},\ldots,\delta\mathbf 1_{[N]})\bigr|\\
 &\qquad\le
 \frac{((1+\eta)\delta)^{k-1}}{NM}
 \sum_{y=1}^M\sum_{x\in\mathbb Z}
 \bigl|(H_j-\delta\mathbf 1_{[N]})(x+P_{j,q}(y))\bigr|\\
 &\qquad=
 \frac{((1+\eta)\delta)^{k-1}}{N}\|H_j-\delta\mathbf 1_{[N]}\|_{\ell^1}
 \le 2\eta(1+\eta)^{k-1}\delta^k.
\end{aligned}
\]
Summing over the \(\ell+1\) terms gives
\[
 \bigl|T_1(H_0,\ldots,H_\ell)-T_1(\delta\mathbf 1_{[N]},\ldots,\delta\mathbf 1_{[N]})\bigr|
 \le 2(\ell+1)\eta(1+\eta)^{k-1}\delta^k.
\] 
For sufficiently small fixed
\(\eta>0\) this is less than \(3c_3\delta^k/8\), which is impossible by \eqref{e:diff1a}. Hence some atom of one
of the factors has density at least \((1+\eta)\delta\), as required.
\end{proof}

\section{Comparison on smooth moduli}\label{s:relative}

We prove Theorem~\ref{t:relative-selector-comparison}, the counting
comparison used for the density increment in
Section~\ref{s:recurrence}. We split the difference of the weights as
\[
 \Lambda_qF_q-\mu_F=(\Lambda_q-1)F_q+(F_q-\mu_F).
\]
For the first term we estimate the prime discrepancy in a full
Gowers norm and control the effect of multiplication by \(F_q\).
For the second term we reduce the count to polynomial-phase
correlations on long progressions and use the arithmetic of the
Heisenberg orbit to bound those correlations.

We begin with the prime discrepancy. On the smooth moduli under
consideration, the Cram\'er model for the normalized weight
\(\Lambda_q\) can be replaced by \(1\).
This gives the estimate for \(\Lambda_q-1\) before inserting the
factor \(F_q\).

\begin{proposition}\label{p:high-order-relative-prime}
Fix integers $d\geq1$ and $S\geq2$. There are
$c_{\rm rel},\eta_{\rm rel}>0$ and an integer $w_0\geq16$,
depending only on $d,S$, such that for all sufficiently large
integers $N$ and all $q,w,M\in\mathbb N$ satisfying
\eqref{e:relative-smooth-hypotheses} and
\begin{equation}\label{e:relative-scale-window}
 N^{1/d}\exp(-d\sqrt{\log N})\leq M\leq N^{1/d},
\end{equation}
one has
\begin{equation}\label{e:high-order-relative-prime}
 \|\Lambda_q-1\|_{U^S[M]}\ll_{d,S}w^{-\eta_{\rm rel}}.
\end{equation}
In particular,
\begin{equation}\label{e:relative-first-moment}
 \mathbb E_{y\leq M}\Lambda_q(y)
 =1+O_{d,S}(w^{-\eta_{\rm rel}}).
\end{equation}
\end{proposition}

\begin{proof}
For a real zero $\beta_0\in(0,1)$ of a primitive nonprincipal real
character of conductor $3\leq r\leq\widetilde Q_N$,
\cite[Theorem~1.1]{BanksSiegel} gives
\[
 (1-\beta_0)\log\widetilde Q_N
 \geq (1-\beta_0)\log r
 \geq 3\cdot10^{-3}.
\]
This excludes the exceptional zero in
\cite[Definition~2.3]{KMTT}, whose threshold is $10^{-3}$.
Thus the model in \cite[Lemma~3.1]{KMTT} is
$\Lambda_{\Cramer,\widetilde Q_N}$.

Take \(c_{\rm rel}\le1/10\) sufficiently small that
\(2c_0c_{\rm rel}<1/10\). With \(Y=qM+1\), the scale conditions give
\(\log Y\asymp_d\log N\), and therefore
\[
 \exp((\log\log Y)^{1/c_S})
 \le\widetilde Q_N\le \exp((\log Y)^{1/10})
\]
for large \(N\), where \(c_S>0\) is the constant of
\cite[Lemma~3.1]{KMTT}. Apply \cite[Lemma~3.1]{KMTT} to the ambient norm at
$Y=qM+1$, then apply the normalized dilation bound
\cite[Lemma~2.2]{KMTT} with length $qM$ and residue $1$. This gives
\[
 \left\|\Lambda_q-\frac{\varphi(q)}q
       \Lambda_{\Cramer,\widetilde Q_N}(q\,\cdot+1)\right\|_{U^S[M]}
 \ll_S q^{(S+1)/2^S}\exp(-(\log\widetilde Q_N)^{c_S}).
\]
After decreasing \(c_{\rm rel}\) so that \(2c_{\rm rel}<c_S\), this
is \(O_{S,d}(w^{-1})\), since \(\log q\le w^2\). The hypotheses of
\cite[Lemma~4.3(i)]{KMTT} are precisely the smoothness and scale
conditions above, with \(b=1\) and \(J=[M]\) in that lemma. It gives
\[
 \left\|\frac{\varphi(q)}q
 \Lambda_{\Cramer,\widetilde Q_N}(q\,\cdot+1)-1\right\|_{U^S[M]}
 \ll_{d,S} w^{-\eta_{\rm rel}}
\]
after taking \(\eta_{\rm rel}\le1\). The triangle inequality proves
\eqref{e:high-order-relative-prime}, and Gowers monotonicity gives \eqref{e:relative-first-moment}.
\end{proof}

To estimate the count weighted by \(F_q-\mu_F\), we use a local
generalized von Neumann estimate. It is uniform over
\((C,q)\)-coefficients at the spatial scale \(N\asymp q^{d-1}M^d\),
and reduces the count to tests against ordinary polynomial phases.

\begin{theorem}[Local generalized von Neumann estimate]\label{t:main}
Fix $C\geq1$ and $d,k\in\mathbb N$, with $k\geq2$. There are
$\eta,A_1,A_2>0$ and $B\in\mathbb N$, depending only on $C,d,k$,
with the following property. Let $N,q\in\mathbb N$ satisfy
$N\geq3$ and $1\leq q\leq\exp((\log N)^{1/10})$, and put
$M=\lfloor(N/q^{d-1})^{1/d}\rfloor\geq1$.
Let $P_1,\ldots,P_{k-1}$ have distinct increasing degrees, largest
degree $d$, and $(C,q)$-coefficients. Suppose that $v$ is $1$-bounded
on $[M]$ and that the $f_i$ are $1$-bounded and supported on
$[-CN,CN]$. For $0<\rho\leq1/2$, one has
\begin{equation}\label{e:main}
 |T_{N,M,v}(\mathbf f)|
 \ll_{C,d,k} N^{-\eta}+\rho+
 \left(\sup_{I,b,u}
 \|v\|_{u^{d+1}[I\cap(u+q^b\mathbb Z)]}\right)^\eta,
\end{equation}
where $b,u\in\mathbb Z$, $0\leq b\leq B$, and $I\subset[M]$ is
an integer interval with $|I|\geq Mq^{-A_1}\rho^{A_2}$ and
$I\cap(u+q^b\mathbb Z)\ne\varnothing$. For $k=2$ one can take
$b=0$ and $I=[M]$.
\end{theorem}

The proof of Theorem~\ref{t:main} is given in
Section~\ref{s:local-vn}.  We state the arithmetic estimate for
those tests using the Heisenberg orbit.

We use the Heisenberg group \(H\) and lattice \(\Gamma\) from the
introduction, with the lower central series filtration. Put
\(g_0=(\alpha,\beta,0)\), so that the orbit \(g_0^n\Gamma\) has vertical
coordinate \(\{Z(n)\}\) in the chart \(0<\{\beta n\}<1\).
A function \(\Phi\) on \(H/\Gamma\) has zero vertical mean if
\(\int_0^1\Phi(g(0,0,z)\Gamma)\,dz=0\) for every \(g\in H\).
We write \(m_H\) for Haar probability measure on \(H/\Gamma\).

\begin{proposition}\label{p:heisenberg-phase}
Fix $d\in\mathbb N$ and $A,L>0$. For sufficiently large $X$,
choose $j$ so that $r_j\leq X^{1/100}<r_{j+1}$.
Suppose $a\in\mathbb Z$ and $v,K\in\mathbb N$, with
$v\leq X^{1/10000}$, $r_j\nmid v$, and $K\geq X^{1/2}$.
For every Lipschitz $\Phi$ on $H/\Gamma$ with
$\|\Phi\|_{\mathrm{Lip}}\leq(\log X)^L$,
\begin{equation}\label{e:heisenberg-equidistribution}
 \left|\mathbb E_{t\leq K}\Phi(g_0^{a+vt}\Gamma)
       -\int_{H/\Gamma}\Phi\,dm_H\right|
 \ll_{A,L,\alpha,\beta}(\log X)^{-A}.
\end{equation}
If $\Phi$ has zero vertical mean, then also
\begin{equation}\label{e:heisenberg-twisted}
 \sup_{\deg\phi\leq d}
 \left|\mathbb E_{t\leq K}
   \Phi(g_0^{a+vt}\Gamma)\e(-\phi(t))\right|
 \ll_{A,L,d,\alpha,\beta}(\log X)^{-A},
\end{equation}
where the supremum is over real polynomials. In particular, for
fixed Lipschitz $F:\mathbb T\to\mathbb C$, with
$\mu_F=\int_{\mathbb T}F$,
\begin{equation}\label{e:heisenberg-centered}
 \sup_{\deg\phi\leq d}
 \left|\mathbb E_{t\leq K}
   \bigl(F(Z(a+vt))-\mu_F\bigr)\e(-\phi(t))\right|
 \ll_{A,d,\alpha,\beta,F}(\log X)^{-A}.
\end{equation}
\end{proposition}

The proof of Proposition~\ref{p:heisenberg-phase} is given in
Section~\ref{s:example}. Combining Theorem~\ref{t:main} with
Proposition~\ref{p:heisenberg-phase} gives the following bound for
the count weighted by \(F_q-\mu_F\).

\begin{lemma}\label{l:unweighted-selector-relative-range}
Fix \(C,d,k\), the parameters of \(Z\), and a Lipschitz function
\(F:\mathbb T\to\mathbb C\). Let \(N\) be sufficiently large and
suppose \(q,w\in\mathbb N\) satisfy \eqref{e:relative-smooth-hypotheses},
where \(w_0\ge16\) and \(c_{\rm rel}>0\) is sufficiently small.
Put \(M=\lfloor(N/q^{d-1})^{1/d}\rfloor\). Let
\(P_1,\ldots,P_{k-1}\) have \((C,q)\)-coefficients and distinct
increasing degrees, with largest degree \(d\). Then for every
fixed \(A>0\) and every tuple \(\mathbf f=(f_0,\ldots,f_{k-1})\)
of \(1\)-bounded functions supported on \([N]\),
\begin{equation}\label{e:bare-selector-comparison}
 |T_{N,M,F_q-\mu_F}(\mathbf f)|\ll_{C,d,k,\alpha,\beta,F,A}w^{-A}.
\end{equation}
\end{lemma}

\begin{proof}
Apply Theorem~\ref{t:main} to \(F_q-\mu_F\), after dividing by a
fixed constant to make it \(1\)-bounded, with \(\rho=w^{-A-1}\).
For sufficiently small \(c_{\rm rel}\), the bound
\(\log q\le w^2\ll(\log N)^{2c_0c_{\rm rel}}\) puts \(q\) in the
range of the theorem. The test progressions \(P\) have step-size \(q^b\)
with \(b=O(1)\), and
\[
 |P|\gg Mq^{-O(1)}w^{-O(1)}=N^{1/d-o(1)}.
\]
Set \(X=qM+1=N^{1/d+o(1)}\). The values \(qy\) for \(y\in P\)
form a progression of length at least \(X^{1/2}\) and step-size
\(q^{b+1}=X^{o(1)}\le X^{1/10000}\) for large \(N\).

By \eqref{e:horizontal-approximants} and \(X^{1/100}<r_{j+1}\),
we have \(\log\log X\le\sqrt{\log r_j}+O(1)\). Thus \(r_j>w\),
since \(w\ll\log X\). As \(q\) is \(w\)-smooth, \(r_j\nmid q^{b+1}\).
The affine parametrization of \(P\) preserves the degree of each
test phase, so \eqref{e:heisenberg-centered} gives
\[
 \|F_q-\mu_F\|_{u^{d+1}[P]}\ll_{F,D}(\log N)^{-D}
\]
for every fixed \(D>0\), uniformly over the test progressions.
Since \(w\ll\log N\), taking \(D\) sufficiently large in
\eqref{e:main} proves \eqref{e:bare-selector-comparison}.
\end{proof}

It remains to control the prime discrepancy with the selector
inserted. Away from the boundary of the Heisenberg chart we
expand \(F_q\) into vertical Fourier modes and use the following
multiplier estimate for each localized mode.

\begin{lemma}[Localized Fourier multipliers]\label{l:vertical-mode-multiplier}
Let \(S\geq3\), let \(I\) be a finite nonempty interval of integers, and let
\(a,v,h\) be integers. Suppose that \(\chi\) is a Lipschitz function
on \(\mathbb T\) supported on an interval of length less than \(1/2\)
inside \((0,1)\). For every function \(f\) on \(I\),
\begin{align}
 \|f(t)\chi(\{\beta(a+vt)\})\e(hZ(a+vt))\|_{U^S[I]}
 &\leq\|\widehat\chi\|_{\ell^1}\|f\|_{U^S[I]}.
 \label{e:vertical-mode-multiplier}
\end{align}
\end{lemma}

The proof of Lemma~\ref{l:vertical-mode-multiplier} is given in
Section~\ref{s:example}. The scalar equidistribution estimate in
Proposition~\ref{p:heisenberg-phase} bounds the proportion of
parameters near the boundary. Together these estimates control
the prime discrepancy after multiplication by \(F_q\), which
completes the comparison.

\begin{proof}[Proof of Theorem~\ref{t:relative-selector-comparison}]
 The difference in
\eqref{e:relative-selector-constant} is the count with weight
\[
 \Lambda_qF_q-\mu_F=(\Lambda_q-1)F_q+(F_q-\mu_F).
\]
Applying Lemma~\ref{l:unweighted-selector-relative-range} with
\(A=1\) gives
\begin{align}\label{e:f-muf}
    \bigl|T_{N,M,F_q-\mu_F}(\mathbf f)\bigr|
 \ll_{C,d,k,\alpha,\beta,F} w^{-1}.
\end{align}
We show that the count
with weight \((\Lambda_q-1)F_q\) is \(O(w^{-\eta_{\rm rel}/2})\).
Let \(0<\varepsilon<1/16\) and choose a smooth cutoff
\(0\le\chi\le1\) supported in \([\varepsilon,1-\varepsilon]\),
equal to \(1\) on \([2\varepsilon,1-2\varepsilon]\), with
\(\|\chi\|_{\mathrm{Lip}}\ll\varepsilon^{-1}\). Use a fixed
partition of unity to write \(\chi=\sum_\nu\chi_\nu\), with each
piece supported on an interval of length less than \(1/2\).
We first split the count according to the cutoff:
\begin{align}\label{e:firstsecondterms}
 T_{N,M,(\Lambda_q-1)F_q}(\mathbf f)
 &=T_{N,M,(\Lambda_q-1)F_q\chi(\beta q\,\cdot)}(\mathbf f)+T_{N,M,(\Lambda_q-1)F_q(1-\chi(\beta q\,\cdot))}
   (\mathbf f).
\end{align}
In the first count, we expand \(F_q\) by the absolutely convergent
Fourier series of \(F\). This gives
\[
\begin{aligned}
 T_{N,M,(\Lambda_q-1)F_q\chi(\beta q\,\cdot)}(\mathbf f)
 &=\sum_\nu\sum_{h}\widehat F(h)\,
   T_{N,M,(\Lambda_q-1)\chi_\nu(\beta q\,\cdot)
     \e(hZ(q\,\cdot))}(\mathbf f).
\end{aligned}
\]
By \cite[Proposition~3.3]{KMTT} and Gowers monotonicity, choose a
common \(S\ge3\) for the weighted count estimate. Applying that
estimate and then Lemma~\ref{l:vertical-mode-multiplier} gives,
for every \(\nu\) and \(h\in\mathbb Z\),
\begin{equation}\label{e:high-order-gvn}
\begin{aligned}
 \bigl|T_{N,M,(\Lambda_q-1)\chi_\nu(\beta q\,\cdot)
                  \e(hZ(q\,\cdot))}(\mathbf f)\bigr|&\ll_{C,d,k,S}
 \|\widehat{\chi_\nu}\|_{\ell^1}\|\Lambda_q-1\|_{U^S[M]}.
\end{aligned}
\end{equation}
The Sobolev estimate for Fourier coefficients gives
\[
 \|\widehat{\chi_\nu}\|_{\ell^1}
 \ll \|\chi_\nu\|_{H^1(\mathbb T)}
 \ll \|\chi_\nu\|_\infty+\|\chi_\nu'\|_2
 \ll \varepsilon^{-1}.
\]
Summing \eqref{e:high-order-gvn} over \(\nu\) and \(h\) gives
\begin{align}\label{e:firstterm}
 \bigl|T_{N,M,(\Lambda_q-1)F_q\chi(\beta q\,\cdot)}(\mathbf f)\bigr|
 &\ll \left(\sum_{h\in\mathbb Z}|\widehat F(h)|\right)
       \left(\sum_\nu\|\widehat{\chi_\nu}\|_{\ell^1}\right)
       \|\Lambda_q-1\|_{U^S[M]}\\
 &\ll_F \varepsilon^{-1}\|\Lambda_q-1\|_{U^S[M]}\nonumber.
\end{align}
Choose \(c_{\rm rel}\) small enough for
Lemma~\ref{l:unweighted-selector-relative-range} and
Proposition~\ref{p:high-order-relative-prime} at order \(S\).
Since \(\log q\le w^2=o(\sqrt{\log N})\), the prescribed \(M\)
satisfies \eqref{e:relative-scale-window}. The proposition gives
\begin{align}\label{e:second2}
    \|\Lambda_q-1\|_{U^S[M]}\ll w^{-\eta_{\rm rel}},
 \qquad 0<\eta_{\rm rel}\le1.
\end{align}
The localized count is therefore \(O_F(\varepsilon^{-1}w^{-\eta_{\rm rel}})\).
For the boundary term, positivity of \(\Lambda_q\) gives
\(|\Lambda_q-1|\le(\Lambda_q-1)+2\). Thus
\begin{align}\label{e:secondterm}
 &\bigl|T_{N,M,(\Lambda_q-1)F_q(1-\chi(\beta q\,\cdot))}
          (\mathbf f)\bigr|\\
 &\le \|F\|_\infty\,
    \mathbb E_{y\le M}|\Lambda_q(y)-1|(1-\chi(\beta qy))\nonumber\\
 &\le \|F\|_\infty\Bigl(
    2\mathbb E_{y\le M}(1-\chi(\beta qy))
    +\bigl|\mathbb E_{y\le M}(\Lambda_q(y)-1)
                        (1-\chi(\beta qy))\bigr|\Bigr)\nonumber\\
 &\ll_F \mathbb E_{y\le M}(1-\chi(\beta qy))
          +\varepsilon^{-1}\|\Lambda_q-1\|_{U^S[M]}\nonumber.
\end{align}
For the last inequality, expand \(1-\chi\) in Fourier modes.
Its Fourier \(\ell^1\) norm is \(O(\varepsilon^{-1})\) by the
Sobolev estimate, and each linear phase preserves the \(U^S\) norm.

Set \(\varepsilon=w^{-\eta_{\rm rel}/2}/32\) and \(X=qM+1\).
For large \(N\), we have \(M\ge X^{1/2}\), \(q\le X^{1/10000}\),
and \(r_j>w\). Since \(q\) is \(w\)-smooth, \(r_j\nmid q\).
We can therefore apply \eqref{e:heisenberg-equidistribution}
to \(1-\chi\) as a function of the horizontal coordinate. This gives
\begin{align}\label{e:first2}
    \mathbb E_{y\le M}(1-\chi(\beta qy))
 =\int_{\mathbb T}(1-\chi)+O((\log X)^{-1})
 \ll\varepsilon.
\end{align}
Here \(\|1-\chi\|_{\mathrm{Lip}}\ll\varepsilon^{-1}\) is bounded
by a power of \(\log X\), the integral is \(O(\varepsilon)\),
and \((\log X)^{-1}\ll\varepsilon\) since \(w\ll\log X\). By \eqref{e:firstsecondterms}, \eqref{e:secondterm}, \eqref{e:second2}, \eqref{e:firstterm}, and \eqref{e:first2}, we conclude that
\[
 \bigl|T_{N,M,(\Lambda_q-1)F_q}(\mathbf f)\bigr|
 \ll_F \varepsilon+\varepsilon^{-1}w^{-\eta_{\rm rel}}
 \ll_F w^{-\eta_{\rm rel}/2}.
\]
Together with the \(O(w^{-1})\) bound for \eqref{e:f-muf}, this proves \eqref{e:relative-selector-constant}
with \(\eta_{\rm sel}=\eta_{\rm rel}/2\).
% Finally, \cite[Proposition~3.3]{KMTT} and Gowers monotonicity give
% \[
%  \bigl|T_{N,M,\Lambda_q}(\mathbf f)-T_{N,M,1}(\mathbf f)\bigr|
%  \ll \|\Lambda_q-1\|_{U^S[M]}
%  \ll w^{-\eta_{\rm rel}}.
% \]
% Combining this with \eqref{e:relative-selector-constant} proves
% \eqref{e:relative-selector-prime}.
\end{proof}

\section{A local generalized von Neumann estimate}\label{s:local-vn}

We prove Theorem~\ref{t:main}, which supplies the reduction to
local phase tests in Lemma~\ref{l:unweighted-selector-relative-range}.
We first use spatial differencing and the inverse theorem for
distinct-degree polynomials to extract linear phases on
progressions for the shifted inputs. We then use those phases to
test the weight on short parameter intervals in fixed residue
classes. 
% The progression lengths in the first step determine the
% length of the parameter intervals in the second.

\begin{lemma}\label{p:localinverse}
Fix \(C\ge1\) and \(d,k\in\mathbb N\), with \(k\ge2\). Let \(N\ge3\),
let \(q\) be a positive integer with
\(q\le\exp((\log N)^{1/10})\), and put
\(M=\lfloor(N/q^{d-1})^{1/d}\rfloor\ge1\).
Let \(P_1,\ldots,P_{k-1}\) have \((C,q)\)-coefficients and distinct
increasing degrees, with largest degree \(d\). Let \(v\) be
\(1\)-bounded on \([M]\), and let
\(\mathbf f=(f_0,\ldots,f_{k-1})\) consist of \(1\)-bounded functions
supported on \([-CN,CN]\).
Suppose that 
\begin{align}\label{e:Tmnvcondition}
    \lvert T_{N,M,v}(\mathbf f)\rvert\ge\delta, \qquad 0<\delta<1/2.
\end{align}
 Then either \(N\ll_{C,d,k}\delta^{-O_{C,d,k}(1)}\)
or the following holds. For a fixed integer \(B=O_d(1)\), put \(L=q^B\).
For each \(1\le i\le k-1\), there is an integer \(V_i\) satisfying
\[
 M^{\deg P_i}(q/\delta)^{-O_{C,d,k}(1)}
 \ll_{C,d,k} V_i\le M^{\deg P_i},
 \qquad LV_i\le N,
\]
and a function \(\Psi_i\) satisfying
\[
 \left|\sum_x f_i(x)\Psi_i(x)\right|
 \gg_{C,d,k}\delta^{O_{C,d,k}(1)}N,
\]
where for every \(j\in\mathbb Z\) and \(a\in\{0,\ldots,L-1\}\), there are
real \(\theta_{i,j,a}\) and \(\gamma_{i,j,a}\) such that
\[
 \Psi_i(x)=\e\bigl(\theta_{i,j,a}x+\gamma_{i,j,a}\bigr)
 \qquad
 \text{for }x\in ((j-1)LV_i,jLV_i]\cap(a+L\mathbb Z).
\]
\end{lemma}

\begin{proof}
Implicit constants depend only on \(C,d,k\). We write
\(\gg\delta^{O(1)}\) for a lower bound by a constant times a
fixed positive power of \(\delta\).

Since \(|v|\leq1\), Cauchy--Schwarz in \(y\) and \eqref{e:Tmnvcondition} give
\[
\begin{aligned}
 \delta^2N^2
 &\leq \mathbb E_{y\leq M}
     \left|\sum_x\prod_{i=0}^{k-1}f_i(x+P_i(y))\right|^2\\
 &=\sum_{h\in\mathbb Z}\mathbb E_{y\leq M}\sum_x
     \prod_{i=0}^{k-1}\Delta_h f_i(x+P_i(y)),
 \qquad
 \Delta_h f(x)=f(x+h)\overline{f(x)}.
\end{aligned}
\]
The equality follows by writing the two spatial variables in the
square as \(x+h\) and \(x\). Since \(P_0=0\), the support of
\(\Delta_h f_0\) restricts the sum to \(O(N)\) values of \(h\),
and each count has modulus \(O(N)\). Taking absolute values in
the sum over \(h\) therefore gives \(\gg\delta^2N\) values of \(h\) such that
$$\left\vert \mathbb E_{y\leq M}\sum_x
     \prod_{i=0}^{k-1}\Delta_h f_i(x+P_i(y))\right\vert \gg \delta^2 N.$$
For each such \(h\), apply \cite[Lemma~3.4]{KMTT}, which uses
\cite[Theorem~2.4]{ShaoWang}, to the functions \(\Delta_h f_i\).
A common translation puts their supports in an interval
\([N']\) with \(N'\asymp_C N\). The small-scale alternative
\(M\ll(q/\delta)^{O(1)}\) gives \(N\ll(q/\delta)^{O(1)}\) by the
definition of \(M\). Since \(q\leq\exp((\log N)^{1/10})\), the
powers of \(q\) can be absorbed into the left-hand side for
large \(N\), giving \(N\ll\delta^{-O(1)}\).
Otherwise, for each \(i\geq1\) the lemma gives a \(1\)-bounded
function \(\varphi_{h,i}\) such that
\begin{align}
 \left|\sum_x\Delta_h f_i(x)\varphi_{h,i}(x)\right|
   &\gg\delta^{O(1)}N,\nonumber\\
 |\varphi_{h,i}(x+q'_h q^{b_h}t)-\varphi_{h,i}(x)|
   &\ll(q/\delta)^{O(1)}M^{-\deg P_i}|t|, \label{e:d9ffbound1}
\end{align}
where \(q'_h\ll\delta^{-O(1)}\) and \(b_h=O_d(1)\). Choose \(B\)
to be the uniform upper bound for \(b_h\) and put \(L=q^B\).
Replacing \(q^{b_h}\) by \(L\) in the second inequality amounts
to replacing \(t\) by \(q^{B-b_h}t\), which only changes the
implicit exponent of \(q/\delta\).

Fix \(i\), and set
\(V_i=\lfloor M^{\deg P_i}(q/\delta)^{-D_0}\rfloor\) for a
sufficiently large fixed \(D_0\geq B\). Outside the alternative
\(N\ll\delta^{-O(1)}\), the quantity inside the floor is at least
\(2\), since \(q=N^{o(1)}\) and \(M=N^{1/d-o(1)}\). This gives
the required lower bound for \(V_i\), while \(D_0\geq B\) gives
\(LV_i\leq M^{\deg P_i}\leq N\).

Intersect the blocks of length \(LV_i\) with the residue classes
modulo \(L\), retaining the full progressions that meet the
support of \(f_i\), including those at the ends. There are
\(O(N/V_i)\) such progressions, each of the form \(P=a+L[V_i]\).

On each \(P\), define \(\widetilde\varphi_{h,i}(a+Lt)\) to be
\(\varphi_{h,i}(a+Lt')\), where \(t'\) is a fixed representative
in \([V_i]\) of the class \(t\bmod q'_h\). Applying \eqref{e:d9ffbound1} with shift \(L(t-t')\) gives
\[
 \bigl|\varphi_{h,i}(a+Lt)-\widetilde\varphi_{h,i}(a+Lt)\bigr|
 \ll (q/\delta)^{O(1)}M^{-\deg P_i}\frac{|t-t'|}{q'_h}
 \ll (q/\delta)^{-D_0+O(1)}.
\]
Since \(\Delta_h f_i\) is \(1\)-bounded and supported on \(O(N)\)
points, the change in the correlation satisfies
\[
 \left|\sum_x\Delta_h f_i(x)
       \bigl(\varphi_{h,i}(x)-\widetilde\varphi_{h,i}(x)\bigr)\right|
 \ll N(q/\delta)^{-D_0+O(1)}.
\]
Choosing \(D_0\) sufficiently large makes this error at most half
the lower bound for the original correlation. It follows that
\[
\begin{aligned}
 \left|\sum_x\Delta_h f_i(x)\widetilde\varphi_{h,i}(x)\right|
 &\ge
 \left|\sum_x\Delta_h f_i(x)\varphi_{h,i}(x)\right|
 -\left|\sum_x\Delta_h f_i(x)
   \bigl(\varphi_{h,i}(x)-\widetilde\varphi_{h,i}(x)\bigr)\right|\\
 &\ge \frac12
 \left|\sum_x\Delta_h f_i(x)\varphi_{h,i}(x)\right|
 \gg \delta^{O(1)}N.
\end{aligned}
\]
On each \(P=a+L[V_i]\), expand the function of \(t\bmod q'_h\) as
\[
 \widetilde\varphi_{h,i}(a+Lt)
 =\sum_{m=0}^{q'_h-1}c_{h,P}(m)\e(mt/q'_h),
 \qquad |c_{h,P}(m)|\le1.
\]
Substituting this into the correlation and taking absolute values
gives
\[
 \delta^{O(1)}N
 \ll \left|\sum_x\Delta_h f_i(x)\widetilde\varphi_{h,i}(x)\right|
 \le \sum_{m=0}^{q'_h-1}\sum_{P=a+L[V_i]}
    \left|\sum_{t\le V_i}\Delta_h f_i(a+Lt)\e(mt/q'_h)\right|.
\]
Thus for each such \(h\) there is a frequency
\(\omega_h=m/q'_h\) for which
\[
 \sum_{P=a+L[V_i]}
    \left|\sum_{t\le V_i}\Delta_h f_i(a+Lt)\e(\omega_h t)\right|
 \gg \frac{\delta^{O(1)}N}{q'_h}
 \gg \delta^{O(1)}N.
\]
Since there are at most \(\le(\max_h q'_h)^2\ll\delta^{-O(1)}\) many $\omega_h$,
one frequency \(\omega\) is chosen for
\(\gg\delta^{O(1)}N\) values of \(h\). For each \(P=a+L[V_i]\), put
\[
 S_P=\sup_{\theta\in\mathbb T}
 \left|\sum_{t\le V_i}f_i(a+Lt)\e(\theta t)\right|:= \left|\sum_{t\le V_i}f_i(a+Lt)\e(\theta_P t)\right|.
\]
The change of variables \(x=a+Lt+h\) gives
\[
 \sum_h\left(\sum_{t\le V_i}\Delta_h f_i(a+Lt)\e(\omega t)\right)
       \e(-h\xi)
 =\widehat f_i(\xi)\sum_{t\le V_i}
   \overline{f_i(a+Lt)}\e\bigl(\omega t+\xi(a+Lt)\bigr).
\]
The sum on the right has modulus at most \(S_P\), by conjugating
and taking \(\theta=-\omega-L\xi\). Plancherel in \(h\) therefore gives
\begin{equation}\label{e:local-plancherel}
 \sum_{h\in\mathbb Z}
 \left|\sum_{t\le V_i}\Delta_h f_i(a+Lt)\e(\omega t)\right|^2
 \le S_P^2\|f_i\|_{\ell^2(\mathbb Z)}^2.
\end{equation}
For each of the \(\gg\delta^{O(1)}N\) shifts with \(\omega_h=\omega\),
the sum over \(P\) is \(\gg\delta^{O(1)}N\). Since
\(\Delta_h f_i=0\) for \(|h|>2CN\), only \(O(N)\) shifts contribute.
There are \(O(N/V_i)\) progressions, so Cauchy--Schwarz in \(h\)
and \(P\), followed by \eqref{e:local-plancherel}, gives
\[
\begin{aligned}
 \delta^{O(1)}N^2
 &\ll \sum_{h:\,\omega_h=\omega}\sum_{P=a+L[V_i]}
      \left|\sum_{t\le V_i}\Delta_h f_i(a+Lt)\e(\omega t)\right|\\
 &\le \sum_{h\in\mathbb Z}\sum_{P=a+L[V_i]}
      \left|\sum_{t\le V_i}\Delta_h f_i(a+Lt)\e(\omega t)\right|\\
 &\ll \frac{N}{\sqrt{V_i}}
      \left(\sum_P\sum_{h\in\mathbb Z}
      \left|\sum_{t\le V_i}\Delta_h f_i(a+Lt)\e(\omega t)\right|^2
      \right)^{1/2}\\
 &\le \frac{N}{\sqrt{V_i}}\,
      \|f_i\|_{\ell^2(\mathbb Z)}
      \left(\sum_P S_P^2\right)^{1/2}.
\end{aligned}
\]
Squaring and using \(\|f_i\|_{\ell^2(\mathbb Z)}^2\ll N\) gives
\(\sum_P S_P^2\gg\delta^{O(1)}NV_i\).
Since \(S_P\leq V_i\), it follows that
\[\sum_P S_P\gg\delta^{O(1)}N.\]
% For each \(P=a+L[V_i]\), choose \(\theta_P\in\mathbb T\) such that
% \[
%  \left|\sum_{t=1}^{V_i} f_i(a+Lt)\e(\theta_Pt)\right|=S_P.
% \]
If \(S_P>0\), choose \(\gamma_P\in\mathbb T\) so that
\[
 \e(\gamma_P)
 =\frac{\overline{\sum_{t=1}^{V_i}f_i(a+Lt)\e(\theta_Pt)}}{S_P}.
\]
If \(S_P=0\), take \(\gamma_P=0\). On each full progression \(P=a+L[V_i]\) meeting the support of
\(f_i\), define
\[
 \Psi_i(x)=\e\!\left(\theta_P\frac{x-a}{L}+\gamma_P\right)
 \qquad (x\in P).
\]
On each progression disjoint from the support of \(f_i\), take
\(\Psi_i(x)=1\), the constant phase.  Thus the contribution
from each \(P\) equals \(S_P\), and
\[
 \sum_x f_i(x)\Psi_i(x)
 =\sum_{P=a+L[V_i]}\sum_{t=1}^{V_i}
      f_i(a+Lt)\e(\theta_Pt+\gamma_P)
 =\sum_P S_P
 \gg \delta^{O(1)}N.
\]
\end{proof}

To finish the reduction, we apply the inverse lemma to the dual
function of each shifted input. This lets us successively replace
the inputs by local linear phases while retaining a large count.
We partition the parameter range into short intervals and discard
the spatial points at which a shift can cross a block boundary.
On each remaining interval and residue class, the shifted arguments
stay in their respective spatial progressions. Evaluating the linear
phases at those shifts gives a single polynomial phase of degree
at most \(d\) in the parameter.

\begin{proof}[Proof of Theorem~\ref{t:main}]
All implicit constants in this proof depend only on $C,d,k$.
For $k=2$, Fourier inversion in $x$ gives
\[
 T_{N,M,v}(f_0,f_1)
 =\frac1N\int_{\mathbb T}\widehat f_0(-\theta)\widehat f_1(\theta)
       \mathbb E_{y\leq M}v(y)\e(\theta P_1(y))\,d\theta.
\]
Cauchy--Schwarz and Plancherel therefore imply
\[
 |T_{N,M,v}(f_0,f_1)|
 \leq \frac{\|f_0\|_{\ell^2}\|f_1\|_{\ell^2}}{N}
       \sup_{\theta\in\mathbb T}
       \left|\mathbb E_{y\leq M}v(y)\e(\theta P_1(y))\right|
 \ll_C \|v\|_{u^{d+1}[M]}.
\]
Now suppose $k\geq3$ and $|T_{N,M,v}(\mathbf f)|\geq\delta$,
where $0<\delta<1/2$.

We replace the shifted inputs one at a time. After each replacement,
the absolute value of the count remains at least a fixed positive
power of \(\delta\), up to a constant depending only on \(C,d,k\). At each stage, $f_j$ denotes the
current input in the $j$th position. For the $i$th replacement, put
\[
 D_i(z)=\mathbb E_{y\leq M}v(y)
       \prod_{\substack{0\leq j\leq k-1\\j\ne i}}
       f_j(z+P_j(y)-P_i(y)).
\]
Changing variables to $z=x+P_i(y)$ gives
$$T_{N,M,v}(\mathbf f)=N^{-1}\sum_z f_i(z)D_i(z).$$
The function $D_i$ is $1$-bounded. Moreover, its support lies in
a fixed enlargement of $[-CN,CN]$, because the product contains
$f_0(z-P_i(y))$ and $\max_{y\leq M}|P_i(y)|\ll N$.
We use this enlarged support interval for every application of
Lemma~\ref{p:localinverse}. 

The phases supplied by Lemma~\ref{p:localinverse} remain
defined on full progressions. When using them as inputs to the
next application, we truncate them to the enlarged interval.
This leaves every count unchanged, since $f_0(x)\ne0$ forces all
arguments $x+P_j(y)$ to lie in that interval.
Since the current count has modulus $\gg\delta^{O(1)}$,
Cauchy--Schwarz gives
\[
 \delta^{O(1)}
 \ll \left|\frac1N\sum_z f_i(z)D_i(z)\right|^2
 \leq \frac{\|f_i\|_{\ell^2}^2}{N^2}\sum_z|D_i(z)|^2
 \ll \frac1N\sum_z|D_i(z)|^2.
\]
Replacing the $i$th input by $\overline{D_i}$ thus gives a
large count, since
\[
 T_{N,M,v}(f_0,\ldots,\overline{D_i},\ldots,f_{k-1})
 =\frac1N\sum_z|D_i(z)|^2
 \gg\delta^{O(1)}.
\]
Unless $N\ll\delta^{-O(1)}$, Lemma~\ref{p:localinverse} supplies
$\Psi_i$ such that
\[
\begin{aligned}
 \left|T_{N,M,v}
       (f_0,\ldots,\overline{\Psi_i},\ldots,f_{k-1})\right|
 &=\frac1N\left|\sum_z\overline{D_i(z)}\Psi_i(z)\right|
 \gg\delta^{O(1)}.
\end{aligned}
\]
This completes the $i$th replacement. Iterating over
$i=1,\ldots,k-1$ gives
\[
 \left|T_{N,M,v}
 \bigl(f_0,\overline{\Psi_1},\ldots,\overline{\Psi_{k-1}}\bigr)\right|
 \gg\delta^{O(1)}.
\] Each $\Psi_i$ is
a linear phase on the intersections of blocks of length $LV_i$
with residue classes modulo the common modulus $L=q^B$.
We next restrict the parameter to intervals on which these block
choices remain constant. Choose a sufficiently large fixed $D_1$
and put $H=\lfloor M(q/\delta)^{-D_1}\rfloor$. If
$M(q/\delta)^{-D_1}<2$, then
\[
 N<q^{d-1}(M+1)^d
 \ll q^{d-1+dD_1}\delta^{-dD_1}.
\]
Since $q\leq\exp((\log N)^{1/10})$, the power of $q$ on the
right is at most $N^{1/2}$ for sufficiently large $N$. This
gives the small-scale alternative $N\ll\delta^{-O(1)}$.
Otherwise,
\[
 H\geq \tfrac12 M(q/\delta)^{-D_1}\geq1.
\]
Partition $[M]$ into intervals of length $H$ and a final interval
of length less than $H$. The final interval contributes $O(H/M)$
to the count.

Let $I$ be a complete interval with first point $y_0$. Writing
$e_i=\deg P_i$, the coefficient bounds give
\[
 |P_i(y)-P_i(y_0)|
 \leq C_1q^{e_i-1}M^{e_i-1}H
 \qquad (y\in I),
\]
where $C_1$ depends only on $C,d$. If
\begin{align}\label{e:distancecondition}
    \operatorname{dist} \bigl(x+P_i(y_0),LV_i\mathbb Z\bigr)
 >C_1q^{e_i-1}M^{e_i-1}H,
\end{align}
then there is an integer $j_i$, depending on $x$ and $I$, such that
\[
 (j_i-1)LV_i<x+P_i(y)\leq j_iLV_i
 \qquad\text{for every }y\in I.
\]
Thus all the points $x+P_i(y)$ lie in the same block as $y$ ranges
over $I$.
There are $O(N/(LV_i))$ relevant boundaries, so the number of
exceptional $x$ for this $I$ and $i$ is at most
$O(Nq^{e_i-1}M^{e_i-1}H/(LV_i))$. Delete these pairs $(x,I)$.
Using the lower bounds for $V_i$, the total change in the
normalized count is
\[
 \ll \frac HM+
       \sum_{i=1}^{k-1}
       \frac{q^{e_i-1}M^{e_i-1}H}{LV_i}
 \ll (q/\delta)^{-D_1+O(1)}.
\]
The count after the replacements has modulus at least
$c\delta^A$, for fixed $c,A>0$. Choose $D_1$ large enough
that the deletion error is at most $\frac12c\delta^A$.
The remaining count then has modulus at least
$\frac12c\delta^A\gg\delta^{O(1)}$.

For each complete interval \(I\), call \(x\) good if it satisfies \eqref{e:distancecondition}
for every \(i\). Summing over these
intervals and their good values of \(x\), we have
\[
\begin{aligned}
 \delta^{O(1)}
 &\ll \Big|
 \frac1{NM}\sum_I\sum_{\substack{x\\\text{good for }I}}
 f_0(x)\sum_{y\in I}v(y)
 \prod_{i=1}^{k-1}\overline{\Psi_i(x+P_i(y))}
 \Big|\\
 &=\Big|
 \sum_I\sum_{\substack{x\\\text{good for }I}}
 \frac{f_0(x)|I|}{NM}
 \mathbb E_{y\in I}v(y)
 \prod_{i=1}^{k-1}\overline{\Psi_i(x+P_i(y))}
 \Big|.
\end{aligned}
\]
The absolute values of the coefficients sum to at most
\[
 \sum_I\sum_{\substack{x\\\text{good for }I}}
 \frac{|f_0(x)||I|}{NM}
 \leq\frac1N\sum_x|f_0(x)|\ll_C1.
\]
Hence there is a complete interval \(I\) and an \(x\) good for
\(I\) such that
\[
 \left|\mathbb E_{y\in I}v(y)
 \prod_{i=1}^{k-1}\overline{\Psi_i(x+P_i(y))}\right|
 \gg\delta^{O(1)}.
\]

On each nonempty class $I\cap(u+L\mathbb Z)$, every
$x+P_i(y)$ stays in a fixed block and a fixed residue class
modulo $L$. On the corresponding spatial progression, write
$\overline{\Psi_i(z)}=\e(\theta_i z+\gamma_i)$.
Their product is $\e(\phi(y))$, where
\[
 \phi(y)=\sum_{i=1}^{k-1}
       \bigl(\theta_i(x+P_i(y))+\gamma_i\bigr),
 \qquad \deg\phi\leq d.
\]
The coefficients may depend on the chosen residue class.
Splitting the average over $I$ into these classes gives
\[
\begin{aligned}
 \delta^{O(1)}
 &\ll
 \left|\mathbb E_{y\in I}v(y)
       \prod_{i=1}^{k-1}\overline{\Psi_i(x+P_i(y))}\right|\leq
 \sum_{\substack{0\leq u<L\\I\cap(u+L\mathbb Z)\ne\varnothing}}
 \frac{|I\cap(u+L\mathbb Z)|}{|I|}
 \|v\|_{u^{d+1}[I\cap(u+L\mathbb Z)]}.
\end{aligned}
\]
The coefficients in this sum are nonnegative and sum to one.
Hence some $u$ satisfies
\[
 \|v\|_{u^{d+1}[I\cap(u+L\mathbb Z)]}\gg\delta^{O(1)},
 \qquad |I|=H\gg Mq^{-D_1}\delta^{D_1}.
\]
This proves the inverse assertion.

To deduce \eqref{e:main}, choose a fixed $D$ large enough that
the small alternative is $N\ll\delta^{-D}$ and, otherwise, the
inverse assertion supplies
\[
 |I|\gg Mq^{-D}\delta^D,
 \qquad
 \|v\|_{u^{d+1}[I\cap(u+L\mathbb Z)]}\gg\delta^D.
\]
Take $\delta$ to be a sufficiently small fixed multiple of
$|T_{N,M,v}(\mathbf f)|$. If $\delta\leq\rho$, the $\rho$ term
in \eqref{e:main} gives the required bound. The small alternative
gives $|T_{N,M,v}(\mathbf f)|\ll N^{-1/D}$. In the remaining case
$\delta>\rho$, taking $A_1,A_2$ sufficiently large ensures that
the interval above has length at least $Mq^{-A_1}\rho^{A_2}$.
It is therefore included in the supremum in \eqref{e:main}, and
\[
 |T_{N,M,v}(\mathbf f)|
 \ll
 \left(\sup_{I,b,u}
 \|v\|_{u^{d+1}[I\cap(u+q^b\mathbb Z)]}\right)^{1/D}.
\]
Choosing $\eta\leq1/D$ proves \eqref{e:main}.
\end{proof}

\section{Phase estimates for the Heisenberg selector}\label{s:example}

We prove Proposition~\ref{p:heisenberg-phase} and
Lemma~\ref{l:vertical-mode-multiplier}, the arithmetic and multiplier
estimates used in Section~\ref{s:relative}. We begin with the horizontal
Diophantine bound needed for equidistribution of the orbit.
We then derive phase cancellation and control multiplication by
localized vertical modes. The same estimates allow us to insert
the von Mangoldt weight for polylogarithmic moduli and obtain the
\(q=1\) input to Theorem~\ref{t:pointwise}.

\begin{lemma}\label{l:heisenberg-horizontal}
For sufficiently large \(X\), choose \(j\) so that
\(r_j\leq X^{1/100}<r_{j+1}\). For every fixed \(L>0\),
\begin{align}\label{e:infcondition}
    \frac{r_j}{(\log X)^L}\longrightarrow\infty
 \qquad\text{as }X\to\infty.
\end{align}
Moreover, for sufficiently large \(X\), if \(v,k,l\) are integers
with \(1\leq v\leq X^{1/5000}\), \(r_j\nmid v\), and
\(0<\max(|k|,|l|)\leq(\log X)^L\), then
\[
 \|v(k\alpha+l\beta)\|_{\mathbb T}\gg_\alpha X^{-1/20}
\]
uniformly in \(v,k,l\).
\end{lemma}

\begin{proof}
Since \(X^{1/100}<r_{j+1}\), estimate \eqref{e:horizontal-approximants}
gives \(\log X\ll\exp(\sqrt{\log r_j})\), which implies  \eqref{e:infcondition}.
We also have $r_{j+1}/{r_j^2}\gg X^{1/200}.$
Indeed, if \(r_j\leq X^{1/400}\), this follows from
\(r_{j+1}>X^{1/100}\). If \(r_j>X^{1/400}\), the growth in
\eqref{e:horizontal-approximants} gives \(r_{j+1}\gg r_j^4\),
so \(r_{j+1}/r_j^2\gg r_j^2>X^{1/200}\).

Suppose \(k\ne0\). Multiplying by \(r_j\) and using
\eqref{e:badly-approximable} for $m=r_j v k$, we obtain
\[
\begin{aligned}
 \|v(k\alpha+l\beta)\|_{\mathbb T}
 &\geq \frac1{r_j}\|r_jvk\alpha\|_{\mathbb T}
       -|l|v\left|\beta-\frac{a_j}{r_j}\right|\geq \frac{c_\alpha}{|k|v r_j^2}
       -O\left(\frac{|l|v}{r_{j+1}}\right).
\end{aligned}
\]
The second term bound follows from the approximation
\(\left|\beta-a_j/r_j\right|\ll 1/r_{j+1}\) in
\eqref{e:horizontal-approximants}. The error is smaller than half the first term for sufficiently
large \(X\), since
\[
 |kl|v^2\frac{r_j^2}{r_{j+1}}
 \ll X^{1/2500-1/200}(\log X)^{2L}=o(1).
\]
Using \(r_j\leq X^{1/100}\), we conclude that
\[
 \|v(k\alpha+l\beta)\|_{\mathbb T}
 \gg_\alpha X^{-1/5000-1/50}(\log X)^{-L}
 \gg X^{-1/20}.
\]

If \(k=0\), then by \eqref{e:infcondition}, for sufficiently large
\(X\) we have \(0<|l|\leq(\log X)^L<r_j\). Since \(r_j\) is prime
and divides none of \(v,l,a_j\), we have
\(\|vla_j/r_j\|_{\mathbb T}\geq 1/r_j\). Thus
\[
\begin{aligned}
 \|vl\beta\|_{\mathbb T}
 &\geq \left\|\frac{vla_j}{r_j}\right\|_{\mathbb T}
       -|l|v\left|\beta-\frac{a_j}{r_j}\right|\geq \frac1{r_j}-O\left(\frac{|l|v}{r_{j+1}}\right).
\end{aligned}
\]
Since
\[
 \frac{|l|v r_j}{r_{j+1}}
 \ll X^{1/5000-1/200}(\log X)^L=o(1),
\]
the error is at most \(1/(2r_j)\) for sufficiently large \(X\).
Consequently, \(\|vl\beta\|_{\mathbb T}\geq 1/(2r_j)
\geq \tfrac12 X^{-1/100}\), which gives the claimed bound.
\end{proof}

The horizontal bound gives equidistribution on long progressions.
To test against a polynomial phase, we apply quantitative
equidistribution to the product of the Heisenberg orbit and a circle
carrying the phase. A possible relation between the two components
can be absorbed into a horizontal twist on a finite cover. The test
function still has zero vertical mean there, which gives cancellation.

\begin{proof}[Proof of Proposition~\ref{p:heisenberg-phase}]
For \(s\in\mathbb N\), let \(D_s(x,y,z)=(sx,sy,s^2z)\).
This automorphism of \(H\) maps \(\Gamma\) into itself and
therefore induces a map on \(H/\Gamma\). We apply its inverse
\(D_{1/s}\) to group elements before passing to the quotient.

We prove \eqref{e:heisenberg-equidistribution} in a form that
also applies to the dilated orbit used below. Fix \(B>0\), and
suppose that there is a Lipschitz function \(\Psi\) on \(H/\Gamma\)
such that
\[
 \left|
 \frac1{K_0}\sum_{t=1}^{K_0}
 \Psi\bigl(D_{1/s}(g_0^{a'+v't})\Gamma\bigr)
 -\int_{H/\Gamma}\Psi
 \right|
 >(\log X)^{-B}\|\Psi\|_{\mathrm{Lip}},
\]
where the integral is with respect to normalized Haar measure.
Here \(s\) is bounded by a fixed power of \(\log X\),
\(1\leq v'\leq X^{1/5000}\), \(r_j\nmid v'\), and
\(K_0\geq X^{1/3}\). The displayed inequality says that the
dilated orbit fails \((\log X)^{-B}\)-equidistribution. Applying the quantitative Leibman theorem
\cite[Theorem~2.9]{GreenTaoNil} with
\(\varepsilon=(\log X)^{-B}\) gives a nonzero horizontal character
\(kx+ly\bmod 1\), with \(\max(|k|,|l|)\leq(\log X)^{O_B(1)}\).
The horizontal coordinates of \(D_{1/s}(g_0^{a'+v't})\) are
\((a'+v't)(\alpha,\beta)/s\), so the character has linear
coefficient \(v'(k\alpha+l\beta)/s\) along the orbit. Since
\(\|c_0+c_1t\|_{C^\infty[K_0]}=K_0\|c_1\|_{\mathbb T}\), the
bound from the theorem reads
\[
 \left\|\frac{(k\alpha+l\beta)(a'+v't)}{s}\right\|_{C^\infty[K_0]}
 =K_0\left\|\frac{v'(k\alpha+l\beta)}{s}\right\|_{\mathbb T}
 \ll_B(\log X)^{O_B(1)}.
\]
By Lemma~\ref{l:heisenberg-horizontal} and
\(\|su\|_{\mathbb T}\leq s\|u\|_{\mathbb T}\), we also have
\[
 K_0\left\|\frac{v'(k\alpha+l\beta)}{s}\right\|_{\mathbb T}
 \geq \frac{K_0}{s}\|v'(k\alpha+l\beta)\|_{\mathbb T}
 \gg_\alpha \frac{X^{1/3-1/20}}{s}.
\]
Since \(s\) is bounded by a fixed power of \(\log X\), this
contradicts the upper bound for sufficiently large \(X\). Thus
the orbit is \((\log X)^{-B}\)-equidistributed under the stated
conditions. Taking \(s=1\), \(a'=a\), \(v'=v\), and \(K_0=K\)
proves \eqref{e:heisenberg-equidistribution}, with \(B>A+L\) to
account for the Lipschitz norm of \(\Phi\).

For \eqref{e:heisenberg-twisted}, suppose the average has modulus
greater than \((\log X)^{-A}\) for some \(\phi\). On
\((H/\Gamma)\times\mathbb T\), equip the circle with its degree-\(d\)
filtration and use the product filtration of degree \(\max(2,d)\).
Consider the polynomial orbit
\((g_0^{a+vt}\Gamma,\phi(t)\bmod 1)\). The test function
\((g\Gamma,z)\mapsto\Phi(g\Gamma)\e(-z)\) has zero Haar integral
and Lipschitz norm \(O((\log X)^L)\). Hence this orbit fails
\((\log X)^{-A-L-1}\)-equidistribution for sufficiently large \(X\).
Quantitative Leibman gives integers \(k,l,m\), not all zero,
of size at most a fixed power of \(\log X\), such that
\[
 \|v(k\alpha+l\beta)t+m\phi(t)\|_{C^\infty[K]}
 \ll(\log X)^{O(1)}.
\]
The logarithmic exponents here and below may depend on \(A,L,d\).
If \(m=0\), then \((k,l)\ne(0,0)\), and
Lemma~\ref{l:heisenberg-horizontal} gives
\[
 K\|v(k\alpha+l\beta)\|_{\mathbb T}
 \gg_\alpha KX^{-1/20}\geq X^{9/20},
\]
contradicting the \(C^\infty[K]\)-bound for large \(X\).
Thus \(m\ne0\).
Write
\[
 v(k\alpha+l\beta)t+m\phi(t)
 =b_0+\sum_{i=1}^d b_i\binom{t}{i}.
\]
For each \(1\leq i\leq d\), choose a nearest integer \(p_i\) to
\(b_i\). The \(C^\infty[K]\)-bound gives
\(\lvert b_i-p_i\rvert=\|b_i\|_{\mathbb T}
\ll K^{-i}(\log X)^{O(1)}\). Set
\[
 P(t)=\sum_{i=1}^d p_i\binom{t}{i},
 \qquad
 E(t)=\sum_{i=1}^d(b_i-p_i)\binom{t}{i}.
\]
Put $c=b_0+a(k\alpha+l\beta)$. Then
\[
 m\phi(t)=-(k\alpha+l\beta)(a+vt)+P(t)+E(t)+c.
\]
To apply equidistribution on the dilated orbit, we decompose the
twist as
\[
 \e(-\phi(t))
 =\e\!\left(\frac{(k\alpha+l\beta)(a+vt)}{m}\right)
  \e(-P(t)/m)\,\e(-E(t)/m)\,\e(-c/m).
\]
The linear factor is absorbed into the test function on the
dilated orbit. The \(P\)-factor is constant on suitable residue
classes, and the \(E\)-factor is handled by partial summation. 

Put \(s=|m|\). The factor \(\e(-P(t)/m)\) is constant on each
residue class modulo \(sd!\). Indeed, \(i!\binom{t}{i}\) is an
integer polynomial in \(t\), so its change under \(t\mapsto t+sd!\)
is divisible by \(sd!\). For \(i\leq d\), division by \(i!\) shows
that \(\binom{t+sd!}{i}-\binom{t}{i}\) is divisible by \(s\).
The same therefore holds for \(P(t+sd!)-P(t)\).

To absorb the linear factor, define
\[
 \widetilde\Phi((x,y,z)\Gamma)
 =\Phi(D_s(x,y,z)\Gamma)
   \e\bigl(\operatorname{sgn}(m)(kx+ly)\bigr).
\]
The horizontal coordinates of the dilated orbit are
\((a+vt)(\alpha,\beta)/s\). Since
\(\operatorname{sgn}(m)/s=1/m\), we obtain
\[
 \Phi(g_0^{a+vt}\Gamma)\e(-\phi(t))
 =\widetilde\Phi(D_{1/s}(g_0^{a+vt})\Gamma)
   \e\bigl(-(P(t)+E(t)+c)/m\bigr).
\]
The Lipschitz norm of \(\widetilde\Phi\) is bounded by a fixed
power of \(\log X\), since \(D_s\) has Lipschitz constant
\(O(s^2)\) in the fixed right-invariant metric. Fix \(x,y\in[0,1)\) and parametrize the vertical fibre by
\((x,y,z)\Gamma\), \(0\leq z<1\). The horizontal phase is
independent of \(z\), so the change of variables \(w=s^2z\)
gives
\[
\begin{aligned}
 \int_0^1\widetilde\Phi((x,y,z)\Gamma)\,dz
 &=\e\bigl(\operatorname{sgn}(m)(kx+ly)\bigr)
   \frac1{s^2}\int_0^{s^2}\Phi((sx,sy,w)\Gamma)\,dw\\
 &=\e\bigl(\operatorname{sgn}(m)(kx+ly)\bigr)
   \int_0^1\Phi((sx,sy,w)\Gamma)\,dw
 =0.
\end{aligned}
\]
The second equality uses \(1\)-periodicity in \(w\) and
\(s^2\in\mathbb N\). The last equality uses the zero vertical
mean of \(\Phi\). Integrating over the horizontal coordinates
now gives \(\int_{H/\Gamma}\widetilde\Phi=0\). 

Fix a residue class modulo \(sd!\), and list its terms in
\([1,K]\) as \(t_1<\cdots<t_R\). The factors \(\e(-P(t_j)/m)\)
and \(\e(-c/m)\) are constant on this class. Removing them and
applying partial summation gives
\begin{align}\label{e:partialsumest}
 \left|\sum_{j=1}^R
   \Phi(g_0^{a+vt_j}\Gamma)\e(-\phi(t_j))\right|
 &=\left|\sum_{j=1}^R
   \widetilde\Phi(D_{1/s}(g_0^{a+vt_j})\Gamma)
   \e(-E(t_j)/m)\right|\\
 &\ll
 \left(1+\sum_{t=1}^{K-1}|E(t+1)-E(t)|\right)
 \max_{1\leq J\leq R}
 \left|\sum_{j=1}^J
   \widetilde\Phi(D_{1/s}(g_0^{a+vt_j})\Gamma)\right|\nonumber.
\end{align}
Indeed, the change in \(\e(-E(t_j)/m)\) between consecutive terms
is bounded by a constant times
\(\sum_{t=t_j}^{t_{j+1}-1}|E(t+1)-E(t)|\). Since
\(\binom{t+1}{i}-\binom{t}{i}=\binom{t}{i-1}\), we have
\begin{align}\label{e:varupestimate}
    \sum_{t=1}^{K-1}|E(t+1)-E(t)|
 \ll (\log X)^{O(1)}\sum_{i=1}^d K^{-i}\binom{K}{i}
 \ll(\log X)^{O(1)}.
\end{align}
We claim that for every fixed \(B>0\),
\[
 \max_{1\leq J\leq R}
 \left|\sum_{j=1}^J
   \widetilde\Phi(D_{1/s}(g_0^{a+vt_j})\Gamma)\right|
 \ll_B R(\log X)^{-B}.
\]
To prove the claim, note that \(a+vt_j=a+v(t_1-sd!)+v'j\),
where \(v'=vsd!\). For sufficiently large \(X\),
\(v'\leq X^{1/10000}(\log X)^{O(1)}\leq X^{1/5000}\).
Also \(r_j\nmid v'\), since \(r_j\nmid v\) and \(sd!<r_j\).
For \(J\geq X^{1/3}\), apply the equidistribution estimate for
the dilated orbit with \(K_0=J\) and exponent \(B+L'\), where
\(L'\) is fixed and
\(\|\widetilde\Phi\|_{\mathrm{Lip}}\leq(\log X)^{L'}\).
Since \(\int_{H/\Gamma}\widetilde\Phi=0\), this gives
\[
 \left|\sum_{j=1}^J
 \widetilde\Phi\bigl(D_{1/s}(g_0^{a+vt_j})\Gamma\bigr)\right|
 \leq J(\log X)^{-B-L'}\|\widetilde\Phi\|_{\mathrm{Lip}}
 \leq J(\log X)^{-B}.
\]
When \(J<X^{1/3}\), the trivial bound is
\(J\|\widetilde\Phi\|_\infty\ll X^{1/3}(\log X)^{O(1)}\).
Since \(R\asymp K/(sd!)\gg X^{1/2}(\log X)^{-O(1)}\), this is
\(O_B(R(\log X)^{-B})\) as well. This proves the claim.

Combining \eqref{e:partialsumest} and \eqref{e:varupestimate}
and taking $B$ sufficiently large gives
\begin{align*}
    \left|\sum_{j=1}^R
   \Phi(g_0^{a+vt_j}\Gamma)\e(-\phi(t_j))\right| \ll_A R(\log X)^{-A-1}.
\end{align*}
Summing over the classes and dividing
by \(K\) gives \(O((\log X)^{-A-1})\) for the twisted average,
contradicting our assumption. This proves
\eqref{e:heisenberg-twisted}.

Right multiplication by $(u,v,w)\in\Gamma$ changes
$z-x\lfloor y\rfloor$ by the integer
$w-u(\lfloor y\rfloor+v)$. Thus $z-x\lfloor y\rfloor\bmod 1$
and $\{y\}$ are well defined on $H/\Gamma$.
Let \(G\) be Lipschitz on \(\mathbb T\) and let \(\chi\) be a
Lipschitz cutoff supported inside \((0,1)\). The function
\[
 (x,y,z)\Gamma\longmapsto
 \chi(\{y\})G(z-x\lfloor y\rfloor)
\]
has Lipschitz norm \(O(\|\chi\|_{\mathrm{Lip}}\|G\|_{\mathrm{Lip}})\).
Indeed, \(z-x\lfloor y\rfloor\) is Lipschitz in bounded local
coordinates away from \(y=0\bmod 1\), and the cutoff makes the
product vanish near that boundary. By translation invariance,
its mean on the vertical fibre over \((x,y)\) is
\[
 \int_0^1\chi(\{y\})G(z-x\lfloor y\rfloor)\,dz
 =\chi(\{y\})\int_{\mathbb T}G.
\]
In particular, taking \(G(u)=\e(hu)\) gives the vertical Fourier
mode of frequency \(h\), with Lipschitz norm
\(O((1+|h|)\|\chi\|_{\mathrm{Lip}})\) and zero vertical mean
when \(h\ne0\).

To prove \eqref{e:heisenberg-centered}, set
\(\varepsilon=(\log X)^{-A-1}\). Choose a smooth cutoff
\(0\leq\chi_\varepsilon\leq1\) that vanishes when
\(\|y\|_{\mathbb T}\leq\varepsilon\), equals \(1\) when
\(\|y\|_{\mathbb T}\geq2\varepsilon\), and has Lipschitz norm
\(O(\varepsilon^{-1})\). Taking \(G=F-\mu_F\) and
\(\chi=\chi_\varepsilon\) gives the test function
\[
 \Phi_\varepsilon((x,y,z)\Gamma)
 =\chi_\varepsilon(\{y\})
  \bigl(F(z-x\lfloor y\rfloor)-\mu_F\bigr).
\]
It has zero vertical mean and Lipschitz norm
\(O_F(\varepsilon^{-1})\).
Evaluating the test function along the orbit gives
\[
 \Phi_\varepsilon(g_0^{a+vt}\Gamma)
 =\chi_\varepsilon(\{\beta(a+vt)\})
  \bigl(F(Z(a+vt))-\mu_F\bigr).
\]
Since $\|\Phi_\varepsilon\|_{\mathrm{Lip}}
\ll_F(\log X)^{A+1}$, applying \eqref{e:heisenberg-twisted} with $L=A+2$ gives
\[
 \left|\mathbb E_{t\leq K}
   \Phi_\varepsilon(g_0^{a+vt}\Gamma)\e(-\phi(t))\right|
 \ll_{A,d,\alpha,\beta,F}(\log X)^{-A}
 \qquad (\deg\phi\leq d).
\]
Removing the cutoff changes the average by at most
\[
 \frac{2\|F\|_\infty}{K}
 \#\bigl\{1\leq t\leq K:
          \|\beta(a+vt)\|_{\mathbb T}\leq2\varepsilon\bigr\}.
\]
A Lipschitz majorant of the strip
\(\|y\|_{\mathbb T}\leq2\varepsilon\), supported where
\(\|y\|_{\mathbb T}\leq3\varepsilon\), has Haar integral
\(O(\varepsilon)\) and Lipschitz norm \(O(\varepsilon^{-1})\).
Applying \eqref{e:heisenberg-equidistribution} to this majorant
with error \(O((\log X)^{-A-1})\) bounds the change by
\(O_F(\varepsilon)\). This proves
\eqref{e:heisenberg-centered}.
\end{proof}



We now prove Lemma~\ref{l:vertical-mode-multiplier}, which
compares the Gowers norms of a localized function before and
after multiplication by a vertical phase.

\begin{proof}[Proof of Lemma~\ref{l:vertical-mode-multiplier}]
We first show that the vertical phase can be removed:
\begin{align}\label{e:twolocalizednorm}
    \|f(t)\chi(\{\beta(a+vt)\})\e(hZ(a+vt))\|_{U^S[I]}
 =\|f(t)\chi(\{\beta(a+vt)\})\|_{U^S[I]}.
\end{align}

Expand the $2^S$-th powers over cubes with vertices
$t_\omega=t+\omega\cdot\mathbf u$. Both expansions have the same
normalizing factor. Zero extension restricts both sums to cubes
with every vertex in $I$, and a zero cutoff at any vertex makes
the corresponding term vanish in both sums.
 On a cube where all cutoffs are nonzero,
write \(n_\omega=a+vt_\omega\) and
\(b_\omega=\lfloor\beta n_\omega\rfloor\). The extra factor in
the cube product on the left is
\[
 \e\!\left(h\sum_{\omega\in\{0,1\}^S}
                (-1)^{|\omega|}Z(n_\omega)\right).
\]
We show that this factor equals \(1\).

On any two-dimensional face, label the four vertices
\(00,01,10,11\), suppressing the fixed coordinates. Since
\(n_{00}+n_{11}=n_{01}+n_{10}\), we have
\[
 b_{00}+b_{11}-b_{01}-b_{10}
 =-\bigl(\{\beta n_{00}\}+\{\beta n_{11}\}
          -\{\beta n_{01}\}-\{\beta n_{10}\}\bigr).
\]
Let \(J\) be the interval supporting \(\chi\). All four
fractional parts lie in \(J\), so
\[
\begin{aligned}
 |b_{00}+b_{11}-b_{01}-b_{10}|
 &=\bigl|\{\beta n_{00}\}+\{\beta n_{11}\}
          -\{\beta n_{01}\}-\{\beta n_{10}\}\bigr|\\
 &\leq \bigl|\{\beta n_{00}\}-\{\beta n_{01}\}\bigr|
       +\bigl|\{\beta n_{11}\}-\{\beta n_{10}\}\bigr|\\
 &\leq 2|J|<1.
\end{aligned}
\]
Since \(b_{00}+b_{11}-b_{01}-b_{10}\) is an integer, it equals zero.
Thus the increment of \(b_\omega\) along an edge in any fixed
direction is independent of the other coordinates. With
\(\mathbf e_i\) denoting the \(i\)th coordinate vector, this gives
\[
 b_\omega=b_{\mathbf 0}
       +\sum_{i=1}^S\omega_i(b_{\mathbf e_i}-b_{\mathbf 0}).
\]
Both \(n_\omega\) and \(b_\omega\) are now affine functions of
\(\omega\). By \eqref{e:heisenberg-phase},
\[
 Z(n_\omega)
 =\frac{\alpha\beta}{2}(n_\omega^2-n_\omega)
  -\alpha n_\omega b_\omega,
\]
Thus \(Z(n_\omega)\) is a polynomial of degree at most two
in \(\omega\). For each of its monomials \(p(\omega)\), choose a
coordinate \(j\) that does not occur in \(p\), which is possible
since \(S\geq3\). Pairing vertices in direction \(j\) gives
\[
 \sum_{\omega\in\{0,1\}^S}(-1)^{|\omega|}p(\omega)
 =\sum_{\substack{\omega\in\{0,1\}^S\\\omega_j=0}}
   (-1)^{|\omega|}
   \bigl(p(\omega)-p(\omega+\mathbf e_j)\bigr)
 =0.
\]
Summing over the monomials, the extra phase factor is
\[
 \e\!\left(h\sum_{\omega\in\{0,1\}^S}
             (-1)^{|\omega|}Z(n_\omega)\right)=\e(0)=1.
\]
Thus the two localized norms in \eqref{e:twolocalizednorm} agree.

To prove \eqref{e:vertical-mode-multiplier}, we expand \(\chi\)
in Fourier modes and bound their contributions separately. To
justify this, we first check that its Fourier coefficients are
absolutely summable. Since \(\chi\) is Lipschitz, its derivative
belongs to \(L^2(\mathbb T)\). Cauchy--Schwarz and Parseval give
\[
 \|\widehat\chi\|_{\ell^1}
 \leq |\widehat\chi(0)|
 +\left(\sum_{k\ne0}\frac1{k^2}\right)^{1/2}
  \left(\sum_{k\ne0}k^2|\widehat\chi(k)|^2\right)^{1/2}
 \ll \|\chi\|_\infty+\|\chi'\|_2<\infty.
\]
Thus the Fourier series converges absolutely and uniformly to
\(\chi\), and
\[
 \chi(\{\beta(a+vt)\})
 =\sum_{k\in\mathbb Z}\widehat\chi(k)\e(k\beta(a+vt)).
\]
Using \eqref{e:twolocalizednorm} and the triangle inequality, we obtain
\[
\begin{aligned}
 \|f(t)\chi(\{\beta(a+vt)\})\e(hZ(a+vt))\|_{U^S[I]}
 &=\|f(t)\chi(\{\beta(a+vt)\})\|_{U^S[I]}\\
 &\leq\sum_{k\in\mathbb Z}|\widehat\chi(k)|
       \|f(t)\e(k\beta(a+vt))\|_{U^S[I]}\\
 &=\|\widehat\chi\|_{\ell^1}\|f\|_{U^S[I]}.
\end{aligned}
\]
The last equality holds because the linear exponent
\(k\beta(a+vt)\) has zero alternating sum on every \(S\)-cube.
\end{proof}


The comparison in Section~\ref{s:relative} replaces
$\Lambda_qF_q$ by the constant $\mu_F$ inside configuration counts
at the smooth moduli of the density increment. There, the high-order
Gowers bound for $\Lambda_q-1$ controls the term $(\Lambda_q-1)F_q$,
and the local generalized von Neumann estimate controls the term
$F_q-\mu_F$ through polynomial-phase tests on progressions.

For pointwise convergence, the polynomial family is fixed.
\cite[Theorem~4.1]{Teravainen} reduces the comparison to bounding
the correlations of $\Lambda(n+1)(F(Z(n))-\mu_F)$ with polynomial
phases on the whole averaging interval. At $q=1$, we retain the
Cram\'er model and use sieve weights to reduce its contribution to
the unweighted phase estimates on progressions. The following
proposition gives the required bound at $q=1$ and also covers
polylogarithmic moduli.

\begin{proposition}[Prime-weighted local phase estimate]
\label{p:prime-phase}
Fix $d\geq1$ and $B,K,L,A>0$. Let $F:\mathbb T\to\mathbb C$
be Lipschitz and let $Z$ be as in \eqref{e:heisenberg-phase}, with
\eqref{e:badly-approximable} and \eqref{e:horizontal-approximants}.
For sufficiently large $N$, suppose that $1\leq q\leq(\log N)^L$
is an integer and put $M=\lfloor(N/q^{d-1})^{1/d}\rfloor$.
If $0\leq b\leq B$ and $u$ are integers, $I\subset[M]$ is an interval
with $|I|\geq M(\log N)^{-K}$, and $P=I\cap(u+q^b\mathbb Z)$, then
\begin{equation}\label{e:prime-local-phase}
 \sup_{\substack{\phi\in\mathbb R[y]\\\deg\phi\leq d}}
 \left|\mathbb E_{y\in P}\Lambda_q(y)
       \bigl(F_q(y)-\mu_F\bigr)\e(-\phi(y))\right|
 \ll_{d,\alpha,\beta,F,A,B,K,L}(\log N)^{-A}.
\end{equation}
\end{proposition}


\begin{proof}
Put \(X=qM+1\). Under the change of variable \(m=qy+1\), the
set \(P\) maps to a progression \(P'\subset[X]\) consisting of
consecutive points of the residue class \(qu+1\) modulo
\(q^{b+1}\). This class is reduced, since \(qu+1\equiv1\pmod q\).
Moreover, \(\log X\asymp_d\log N\), and
\[
 |P'|=|P|\geq \frac{|I|}{q^b}-1
 \gg \frac{M}{q^b}(\log N)^{-K}
 \gg X(\log X)^{-K-(B+1)L}.
\]
Since \(\varphi(q)/q\leq1\), it suffices to bound
\[
 \left|\mathbb E_{m\in P'}\Lambda(m)
 \bigl(F(Z(m-1))-\mu_F\bigr)\e(-\phi((m-1)/q))\right|.
\]
The polynomial \(\phi((m-1)/q)\) has degree at most \(d\).

\medskip
\noindent\emph{Reduction to localized Fourier modes.}
Set \(\eta=(\log X)^{-A-2}\). We claim that it suffices to prove,
for every fixed \(A'>0\),
\begin{align}\label{e:reductiontoexpfun}
    \left|\mathbb E_{m\in P'}\Lambda(m)g(m)\right|
 &\ll_{A'}(\log X)^{-A'}\\
 g(m)&=\chi(\{\beta(m-1)\})
       \e\bigl(hZ(m-1)-\phi((m-1)/q)\bigr), \nonumber
\end{align}
uniformly for \(0<|h|\leq\eta^{-2}\) and for the finitely many
cutoffs \(\chi\) constructed below.

For $0<\delta<1/8$, choose a smooth cutoff
$\chi_\delta:\mathbb T\to[0,1]$ with Lipschitz norm
$O(\delta^{-1})$, equal to zero on $\|y\|_{\mathbb T}\leq\delta$
and to one on $\|y\|_{\mathbb T}\geq2\delta$.

Fix a smooth partition of unity $\psi_0,\ldots,\psi_4$ on $[0,1]$
subordinate to intervals of length less than $1/2$, and set
$\chi_\nu=\chi_\eta\psi_\nu$. Then
$\chi_\eta=\sum_{\nu=0}^4\chi_\nu$. Since $\chi_\eta$ vanishes
near the endpoints, each $\chi_\nu$ extends smoothly to
$\mathbb T$, is supported in an interval of length less than
$1/2$ inside $(0,1)$, and satisfies
$0\leq\chi_\nu\leq1$ and
$\|\chi_\nu\|_{\mathrm{Lip}}\ll\eta^{-1}$.
Since \(0\leq 1-\chi_\eta(y)\leq
\mathbf1_{\|y\|_{\mathbb T}\leq2\eta}\) and
\(\Lambda(m)\leq\log X\) on \([X]\), we have
\[
\begin{aligned}
 &\mathbb E_{m\in P'}\Lambda(m)
   \bigl(F(Z(m-1))-\mu_F\bigr)\e(-\phi((m-1)/q))\\
 &\quad=
 \mathbb E_{m\in P'}\Lambda(m)\chi_\eta(\{\beta(m-1)\})
   \bigl(F(Z(m-1))-\mu_F\bigr)\e(-\phi((m-1)/q))\\
 &\qquad+
 O_F\!\left(
 \log X\,\mathbb E_{m\in P'}
 \mathbf1_{\|\beta(m-1)\|_{\mathbb T}\leq2\eta}
 \right).
\end{aligned}
\]
Write \(m-1=a+q^{b+1}t\) for \(m\in P'\), with
\(1\leq t\leq |P'|\). Choose $j$ at scale $X$ as in
Proposition~\ref{p:heisenberg-phase}. The lower bound for $|P'|$
gives $|P'|\geq X^{1/2}$. The step $q^{b+1}$ is bounded by a
fixed power of $\log X$, so for large $X$ it is at most
$X^{1/10000}$ and is smaller than $r_j$ by
\eqref{e:infcondition}. The second horizontal coordinate of
\(g_0^{m-1}\Gamma\) is \(\{\beta(m-1)\}\). We apply
\eqref{e:heisenberg-equidistribution} to
\(\Phi((x,y,z)\Gamma)=1-\chi_{2\eta}(\{y\})\), whose Lipschitz
norm is \(O(\eta^{-1})\). Haar probability measure on
\(H/\Gamma\) projects to uniform measure on the second
horizontal coordinate, so
\[
\begin{aligned}
 \mathbb E_{m\in P'}
 \mathbf1_{\|\beta(m-1)\|_{\mathbb T}\leq2\eta}
 &\leq \mathbb E_{m\in P'}\Phi(g_0^{m-1}\Gamma)\\
 &=\int_{H/\Gamma}\Phi\,dm_H+O(\eta)
 =\int_{\mathbb T}(1-\chi_{2\eta}(y))\,dy+O(\eta)
 \ll\eta.
\end{aligned}
\]
Here \(1-\chi_{2\eta}\) equals one on
\(\|y\|_{\mathbb T}\leq2\eta\) and is supported on
\(\|y\|_{\mathbb T}\leq4\eta\). Since \(\eta\) is a fixed
negative power of \(\log X\), the error in
\eqref{e:heisenberg-equidistribution} can be taken to be
\(O(\eta)\). Thus inserting the cutoff changes the original
average by \(O_F(\eta\log X)\).

We now truncate the Fourier series of \(F-\mu_F\).
Cauchy--Schwarz and Parseval give
\[
 \sum_{|h|>\eta^{-2}}|\widehat F(h)|
 \ll \|F'\|_2
       \left(\sum_{|h|>\eta^{-2}}h^{-2}\right)^{1/2}
 \ll_F\eta,
 \qquad
 \sum_{h\in\mathbb Z}|\widehat F(h)|\ll_F1.
\]
Thus the truncation changes the average by \(O_F(\eta\log X)\),
using \(\Lambda(m)\leq\log X\). Since \(\widehat F(0)=\mu_F\) and
\(\chi_\eta=\sum_{\nu=0}^4\chi_\nu\), the truncated average is
\[
 \sum_{\nu=0}^4\sum_{0<|h|\leq\eta^{-2}}\widehat F(h)\,
 \mathbb E_{m\in P'}\Lambda(m)\chi_\nu(\{\beta(m-1)\})
 \e\bigl(hZ(m-1)-\phi((m-1)/q)\bigr).
\]
The sum of the absolute values of the coefficients is bounded by
\[
 \sum_{\nu=0}^4\sum_{0<|h|\leq\eta^{-2}}|\widehat F(h)|
 \leq 5\|\widehat F\|_{\ell^1}\ll_F 1.
\]
It therefore suffices to bound each localized average by
\(O((\log X)^{-A})\), uniformly in \(\nu,h,\phi\). The cutoff
and truncation errors are \(O_F(\eta\log X)\).

\medskip
\noindent\emph{The difference between \(\Lambda\) and the Cram\'er model.}
Take \(S=\max(3,d+1)\) and let
\(\widetilde Q_X=\lceil\exp((\log X)^{c_0})\rceil\), with \(c_0\)
as in \cite[Definition~2.3]{KMTT}. We claim that
\begin{align}\label{e:diffgowersnorm}
     \left|\mathbb E_{m\in P'}
   (\Lambda-\Lambda_{\Cramer,\widetilde Q_X})(m)g(m)\right|
 \ll_S
 \frac{X}{|P'|}\log X\,\|\widehat\chi\|_{\ell^1}
 \|\Lambda-\Lambda_{\Cramer,\widetilde Q_X}\|_{U^S[X]}.
\end{align}
To prove this, write \(\mathbf1_{P'}\) as the product of the
indicator of the interval spanned by \(P'\) and the residue-class
indicator
\[
 \mathbf1_{m\equiv qu+1\!\!\!\pmod{q^{b+1}}}
 =\frac1{q^{b+1}}\sum_{r=0}^{q^{b+1}-1}
   \e\!\left(\frac{r(m-qu-1)}{q^{b+1}}\right).
\]
The coefficients in this expansion have absolute sum one. The
Fourier coefficients of the interval indicator on the cyclic
group of order \(2X+1\) have absolute sum \(O(\log X)\), by the
geometric-series bound. Expanding both indicators therefore bounds
the left-hand side by
\[
 \ll \frac{X}{|P'|}\log X
 \sup_{\theta\in\mathbb T}
 \left|\mathbb E_{m\leq X}
   (\Lambda-\Lambda_{\Cramer,\widetilde Q_X})(m)
   g(m)\e(\theta m)\right|.
\]
By Gowers monotonicity, each average in the supremum is bounded
by a constant times the \(U^S[X]\)-norm of its integrand.
The exponent \(\theta m-\phi((m-1)/q)\) has degree at most \(d<S\),
so multiplication by its exponential preserves this norm.
Applying \eqref{e:vertical-mode-multiplier} with \(a=-1\) and
\(v=1\) to the remaining factors gives \eqref{e:diffgowersnorm}.

Since \(\|\widehat\chi\|_{\ell^1}\ll
\|\chi\|_\infty+\|\chi'\|_2\ll\eta^{-1}\) and \(X/|P'|\) is
bounded by a fixed power of \(\log X\), it remains to bound the
ambient Gowers norm. We claim that
\begin{equation}\label{e:ambient-cramer}
 \|\Lambda-\Lambda_{\Cramer,\widetilde Q_X}\|_{U^S[X]}
 \ll_S\exp\bigl(-c(\log X)^{c_0c_S}\bigr)
\end{equation}
for some \(c>0\), where \(c_S\) is the constant in
\cite[Lemma~3.1]{KMTT}.

Indeed, for any real zero \(\beta_\chi\in(0,1)\) of a primitive
nonprincipal real character of conductor \(Q\leq\widetilde Q_X\),
\cite[Theorem~1.1]{BanksSiegel} gives
\[
 (1-\beta_\chi)\log\widetilde Q_X
 \geq (1-\beta_\chi)\log Q
 \geq 3\cdot10^{-3}.
\]
This excludes the exceptional zero in \cite[Definition~2.3]{KMTT},
which would require the left-hand side to be at most \(10^{-3}\).
Thus the Siegel model at this cutoff is the Cram\'er model.
Moreover, for sufficiently large \(X\), the choice \(0<c_0<1/11\)
gives
\[
 \exp((\log\log X)^{1/c_S})
 \leq \widetilde Q_X
 \leq \exp((\log X)^{1/10}).
\]
Taking both scales in \cite[Lemma~3.1]{KMTT} to be \(X\) now
proves \eqref{e:ambient-cramer}.
Substituting \eqref{e:ambient-cramer} into the localized bound \eqref{e:diffgowersnorm}
gives \(O_{A'}((\log X)^{-A'})\) for every fixed \(A'>0\),
uniformly in the localized modes and the polynomial test.

Using \eqref{e:reductiontoexpfun} and \eqref{e:diffgowersnorm}, it remains to prove that, for every fixed \(A'>0\),
\begin{align}\label{e:cramerreduest}
     \left|\mathbb E_{m\in P'}
   \Lambda_{\Cramer,\widetilde Q_X}(m)g(m)\right|
 \ll_{A'}(\log X)^{-A'},
\end{align}
uniformly in \(\chi,h,\phi\). 

\medskip
\noindent\emph{The Cram\'er term.}
Let $W$ be the product of the primes $p\leq\widetilde Q_X$ that
do not divide $q$. Since $m\equiv1\pmod q$ for every $m\in P'$,
\[
 \mathbb E_{m\in P'}
 \Lambda_{\Cramer,\widetilde Q_X}(m)g(m)
 =
 \prod_{p\leq\widetilde Q_X}(1-1/p)^{-1}
 \mathbb E_{m\in P'}\mathbf1_{(m,W)=1}g(m).
\]
The product is $O(\log X)$, so it remains to prove
\[
 \left|\mathbb E_{m\in P'}\mathbf1_{(m,W)=1}g(m)\right|
 \ll_{A'}(\log X)^{-A'-1}.
\]

We approximate the coprimality indicator by a sum of divisibility
indicators with $\ell\leq\mathcal D=X^{1/20000}$. Put $H=|P'|$
and let $H_\ell$ count the $m\in P'$ with $\ell\mid m$. For
$\ell\mid W$, we have $(\ell,q^{b+1})=1$, so the Chinese
remainder theorem puts these $m$ in one residue class modulo
$q^{b+1}\ell$ within the interval spanned by $P'$.  Hence $H_\ell=H/\ell+O(1)$. We apply the fundamental lemma of the sieve
\cite[Lemma~6.3]{IwaniecKowalski} with prime cutoff
$\widetilde Q_X+1$, level $\mathcal D$, and multiplicative local
density given on primes by $1/p$ for $p\mid W$ and zero otherwise.
For \(2\leq z_1<z_2\leq \mathcal D\), Mertens' product
formula gives
\[
 \prod_{\substack{z_1\leq p<z_2\\p\mid W}}
       \left(1-\frac1p\right)^{-1}
 \leq
 \prod_{z_1\leq p<z_2}\left(1-\frac1p\right)^{-1}
 \ll \frac{\log z_2}{\log z_1}.
\]
The first inequality holds because omitting the primes dividing
\(q\) removes factors greater than one.  The sieve parameter is
\(\log\mathcal D/\log(\widetilde Q_X+1)
\asymp(\log X)^{1-c_0}\). Hence the lemma supplies coefficients
\(\lambda_\ell^\pm\) with \(|\lambda_\ell^\pm|\leq1\), supported
on \(\ell\mid W\) and \(\ell\leq\mathcal D\), such that
\[
 \sum_{\ell\mid(m,W)}\lambda_\ell^-
 \leq\mathbf1_{(m,W)=1}
 \leq\sum_{\ell\mid(m,W)}\lambda_\ell^+,
 \qquad
 \sum_{\ell\mid W}\frac{\lambda_\ell^+-\lambda_\ell^-}{\ell}
 \ll_{A'}(\log X)^{-A'-2}.
\]
The midpoint of the two sieve sums thus has average absolute
error at most
\[
\begin{aligned}
 &\mathbb E_{m\in P'}
 \left|\mathbf1_{(m,W)=1}
       -\sum_{\ell\mid(m,W)}
         \frac{\lambda_\ell^++\lambda_\ell^-}{2}\right|\\
 &\qquad\leq
 \frac12\sum_{\ell\mid W}
    (\lambda_\ell^+-\lambda_\ell^-)\frac{H_\ell}{H}
 \ll_{A'}(\log X)^{-A'-2}+\frac{\mathcal D}{H}.
\end{aligned}
\]
The last bound uses $H_\ell/H=1/\ell+O(1/H)$ and the support
of the coefficients. Since $|g|\leq1$ and the midpoint
coefficients have modulus at most one, we obtain
\[
 \left|\mathbb E_{m\in P'}\mathbf1_{(m,W)=1}g(m)\right|
 \leq
 \sum_{\substack{\ell\mid W\\\ell\leq\mathcal D}}
 \left|\mathbb E_{m\in P'}\mathbf1_{\ell\mid m}g(m)\right|
 +O_{A'}\!\left((\log X)^{-A'-2}+\frac{\mathcal D}{H}\right).
\]
Since $\mathcal D/H$ is smaller than every fixed inverse power
of $\log X$,  it remains to show that
\[
 \sum_{\substack{\ell\mid W\\\ell\leq\mathcal D}}
 \left|\mathbb E_{m\in P'}\mathbf1_{\ell\mid m}g(m)\right|
 \ll_{A'}(\log X)^{-A'-1}.
\]

For these divisors, the progression step satisfies
\(q^{b+1}\ell\leq X^{1/10000}\), and \(H_\ell\geq X^{1/2}\)
for large \(X\), by the lower bound for \(H\). If \(r_j\nmid\ell\),
then \(r_j\nmid q^{b+1}\ell\). Write the progression as
\(m-1=a+q^{b+1}\ell t\), \(1\leq t\leq H_\ell\).
The mode with cutoff \(\chi\) has zero vertical mean, since
\(h\ne0\), and its Lipschitz norm is
\(O((1+|h|)\|\chi\|_{\mathrm{Lip}})=O(\eta^{-3})\).
The polynomial test
\(\phi((a+q^{b+1}\ell t)/q)\) has degree at most \(d\).
Apply \eqref{e:heisenberg-twisted} with saving exponent $A'+2$.
For $r_j\nmid\ell$ this gives
\[
 \left|\mathbb E_{m\in P'}\mathbf1_{\ell\mid m}g(m)\right|
 \ll_{A'}\frac{H_\ell}{H}(\log X)^{-A'-2}.
\]
For $r_j\mid\ell$, use the bound $H_\ell/H$. Since
$H_\ell/H\ll1/\ell$ for $\ell\leq\mathcal D$, the sum over
the divisors is
\[
 \ll_{A'}(\log X)^{-A'-2}\sum_{\ell\leq\mathcal D}\frac1\ell
 +\sum_{\substack{\ell\leq\mathcal D\\r_j\mid\ell}}\frac1\ell
 \ll_{A'}(\log X)^{-A'-1}+\frac{\log X}{r_j}.
\]
Both $(\log X)/r_j$ and $\mathcal D/H$ are smaller than any
fixed inverse power of $\log X$. The Cram\'er normalization
costs one further factor of $O(\log X)$, proving
\eqref{e:cramerreduest}. Combining the two localized estimates
with the cutoff and Fourier truncation errors
$O_F(\eta\log X)$ proves \eqref{e:prime-local-phase}.
% Combining the two prime estimates bounds every localized
% average by an arbitrary inverse power of \(\log X\).
% The cutoff and Fourier truncation errors are
% \(O_F(\eta\log X)\), and \(\log X\asymp_d\log N\).
% This proves \eqref{e:prime-local-phase}.
\end{proof}

% \subsection{Comparison for polylogarithmic moduli}

% We now compare the count with the selector to \(\mu_F\) times
% the unrestricted prime-weighted count. The difference is the
% count with weight \(\Lambda_q(F_q-\mu_F)\). When \(q\) is
% polylogarithmic, the local generalized von Neumann theorem asks
% for phase tests on progressions of length at least \(M\) divided
% by a fixed power of \(\log N\). The prime-weighted phase estimate
% above controls these tests, with enough saving to absorb the
% normalization of the weight by \(O(\log N)\).

% \begin{theorem}[Uniform prime--nil-Bohr comparison]\label{t:comparison}
% Fix $C,d,\ell$, $L,A>0$, and a Lipschitz function $F:\mathbb T\to\mathbb C$.
% Let $Z$ be given by \eqref{e:heisenberg-phase} under
% \eqref{e:badly-approximable} and \eqref{e:horizontal-approximants}.
% For sufficiently large $N$, let $1\leq q\leq(\log N)^L$ and
% $M=\lfloor(N/q^{d-1})^{1/d}\rfloor$. If the $R_i$ have distinct
% increasing degrees, largest degree $d$, and $(C,q)$-coefficients,
% then for all $1$-bounded $f_i$ supported on $[-CN,CN]$,
% \[
%  T^{\mathbf R}_{N,M,\Lambda_q F_q}(\mathbf f)
%  =\mu_F T^{\mathbf R}_{N,M,\Lambda_q}(\mathbf f)
%  +O_{C,d,\ell,\alpha,\beta,L,F,A}((\log N)^{-A}).
% \]
% The constants are uniform over the stated coefficient class.
% \end{theorem}

% \begin{proof}
% The difference is the count with weight $\Lambda_q(F_q-\mu_F)$.
% Since $\Lambda_q(y)\leq\log(qM+1)$ on $[M]$, dividing this weight
% by a sufficiently large constant times $\log N$ makes it $1$-bounded.
% Apply Theorem~\ref{t:main} to the normalized weight with
% $\rho=(\log N)^{-A-2}$. The intervals in its phase tests have length at
% least $M(\log N)^{-K}$ with $K=A_1L+A_2(A+2)$, in the notation of
% that theorem. Proposition~\ref{p:prime-phase} bounds all the resulting
% phase norms by $O((\log N)^{-D})$ for any fixed $D>0$. Take
% $D\eta>A+2$ and multiply the bound in \eqref{e:main} by the
% normalizing factor.
% \end{proof}

\section{Pointwise ergodic averages}\label{s:ergodic}

We prove Theorem~\ref{t:pointwise} by comparing the selected and
unrestricted averages. Proposition~\ref{p:prime-phase} at \(q=1\)
gives cancellation of the centered weight
\(\Lambda(n+1)(F(Z(n))-\mu_F)\) against polynomial phases.
Since the polynomial family is fixed, %Ter\"av\"ainen's
second author's
fixed-family estimate converts this cancellation into an
\(L^1\) bound for the difference of the ergodic averages.

We choose bounded measurable representatives of the $f_i$, with
$|f_i(x)|\leq\|f_i\|_\infty$ everywhere. Since $T$ preserves measure,
these choices do not affect any of the averages outside the union
of the integer translates of the exceptional null sets.

For a Lipschitz function \(F:\mathbb T\to\mathbb C\), put
\[
 \mathcal A_H^F(\mathbf f)(x)
 :=\frac1H\sum_{n\le H}\Lambda(n+1)F(Z(n))
       \prod_{i=1}^{\ell} f_i(T^{P_i(n)}x).
\]
Theorem~\ref{t:pointwise} follows from the following comparison for
Lipschitz weights.

\begin{proposition}\label{p:ergodic-comparison}
Let $P_1,\ldots,P_\ell\in\mathbb Z[y]$ have zero constant terms
and distinct positive degrees. Let $\alpha$ be badly approximable
and let $\beta\in(0,1)$ be the limit of an increasing sequence
of rationals $a_j/r_j$ with prime denominators $r_j\to\infty$
satisfying
\[
 r_{j+1}\asymp\exp\!\bigl(\exp(\sqrt{\log r_j})\bigr),
 \qquad 0<\beta-a_j/r_j\ll r_{j+1}^{-1}.
\]
Set $Z(n)=\alpha\beta n(n-1)/2-\alpha n\lfloor\beta n\rfloor$.
Let $(X,\nu,T)$ be an invertible probability measure-preserving
system, let $f_1,\ldots,f_\ell\in L^\infty(X)$, and fix a Lipschitz
function $F:\mathbb T\to\mathbb C$ with mean
$\mu_F=\int_{\mathbb T}F$. Then for every $A>0$ and every integer
$H\geq2$,
\begin{equation}\label{e:ergodic-L1-log}
 \|\mathcal A_H^F(\mathbf f)-\mu_F\mathcal A_H^1(\mathbf f)\|_{L^1(X)}
 \ll_{\mathbf P,\alpha,\beta,F,A}
       (\log H)^{-A}\prod_{i=1}^{\ell}\|f_i\|_\infty.
\end{equation}
The difference tends to zero almost everywhere. Moreover,
\begin{equation}\label{e:smooth-prime-comparison}
 \frac1{\pi(H)}\sum_{p\le H}
 (F(Z(p-1))-\mu_F)\prod_{i=1}^{\ell} f_i(T^{P_i(p-1)}x)
 \longrightarrow 0
 \qquad\text{for almost every }x.
\end{equation}
\end{proposition}

To deduce the theorem, we use the comparison with Lipschitz
weights in two ways. For the given inputs it controls the average
with a Lipschitz minorant of the selector. With constant inputs
it controls the contribution of the gap between the minorant and
a majorant, and gives the asymptotic number of selected primes.

\begin{proof}[Proof of Theorem~\ref{t:pointwise} assuming
Proposition~\ref{p:ergodic-comparison}]
The polynomials \(P_i(t-1)\) have distinct degrees, and the
bounded functions \(f_i\) belong to \(L^{2\ell}(X)\). By
\cite[Theorem~1.1(ii)]{MirekWanWright}, the limit
\[
 L(x):=\lim_{N\to\infty}
 \frac1N\sum_{n\leq N}\Lambda(n)
 \prod_{i=1}^{\ell} f_i(T^{P_i(n-1)}x)
\]
exists for almost every \(x\). Removing higher prime powers and
using the prime number theorem and partial summation gives
\[
 A_N^T(\mathbf f)(x)\longrightarrow L(x)
 \qquad\text{for almost every }x,
\]
where \(A_N^T(\mathbf f)\) is defined in
\eqref{e:prime-ergodic-average}.

For each \(j\ge1\), choose Lipschitz functions \(F_j^-\) and \(F_j^+\)
on \(\mathbb T\) satisfying \(0\le F_j^-\le\mathbf1_J\le F_j^+\le1\)
and \(\int_{\mathbb T}(F_j^+-F_j^-)<2^{-j}\). Apply
\eqref{e:smooth-prime-comparison} to \(F_j^-\) with the given
functions, and to \(F_j^+-F_j^-\) with all functions equal to one.
For almost every $x$, we have $A_N^T(\mathbf f)(x)\to L(x)$
and both comparisons hold for every $j\geq1$. For such $x$,
the inequality $0\leq\mathbf1_J-F_j^-\leq F_j^+-F_j^-$ gives
\[
\begin{aligned}
 &\left|\frac1{\pi(H)}\sum_{p\leq H}
   (\mathbf1_J-F_j^-)(Z(p-1))
   \prod_{i=1}^{\ell}f_i(T^{P_i(p-1)}x)\right|\\
 &\qquad\leq
 \left(\prod_{i=1}^{\ell}\|f_i\|_\infty\right)
 \frac1{\pi(H)}\sum_{p\leq H}(F_j^+-F_j^-)(Z(p-1)).
\end{aligned}
\]
By \eqref{e:smooth-prime-comparison} with all inputs equal to one,
the prime average on the right converges to
\(\int_{\mathbb T}(F_j^+-F_j^-)<2^{-j}\). Hence the limit superior
of the left-hand side is at most
\(2^{-j}\prod_{i=1}^{\ell}\|f_i\|_\infty\).
Since \(|\mu_{F_j^-}-|J||<2^{-j}\), letting \(j\to\infty\) gives
\[
 \frac1{\pi(H)}
 \sum_{\substack{p\le H\\p\in\mathcal P_{Z,J}}}
       \prod_{i=1}^{\ell} f_i(T^{P_i(p-1)}x)
 =|J| A_H^T(\mathbf f)(x)+o(1)
\]
almost everywhere. Taking every \(f_i=1\) gives
\(\pi_{Z,J}(H)\sim |J|\pi(H)\). Division by
\(\pi_{Z,J}(H)/\pi(H)\) completes the proof.
\end{proof}

\begin{proof}[Proof of Proposition~\ref{p:ergodic-comparison}]
By homogeneity assume that the \(f_i\) are \(1\)-bounded. Relabel
the polynomials in increasing order of degree and permute the
corresponding inputs together with them. The products in the
averages and of the input norms are unchanged. Write
$$  \qquad \theta_H(n):=\Lambda(n+1)(F(Z(n))-\mu_F), \quad n\in [H].$$
Proposition~\ref{p:prime-phase}, with \(q=1\), \(N=H^d\), and
the full interval \([H]\), gives
\[\|\theta_H\|_{u^{d+1}[H]}\ll_{\mathbf P,\alpha,\beta,F,B}(\log H)^{-B}, \quad B>0.\] 
Put
\(L_H=2(1+\|F\|_\infty)\log(H+2)\), so that \(\theta_H/L_H\)
is \(1\)-bounded.

Let \(g:X\to\mathbb C\) be \(1\)-bounded and choose \(C\) depending
on the fixed polynomials so that \([H^d]+P_i([H])\subset[-CH^d,CH^d]\)
for all \(i\). For \(z\in X\) and \(i\ge1\), set
\[g_{0,z}(m):=g(T^m z)\mathbf1_{[H^d]}(m),\qquad g_{i,z}(m):=f_i(T^m z)\mathbf1_{[-CH^d,CH^d]}(m).\] 
Averaging the measure-preserving change of variables \(x=T^m z\) over \(m\in[H^d]\) gives
\begin{align*}
 &\int_X (\mathcal A_H^F(\mathbf f)-\mu_F\mathcal A_H^1(\mathbf f))g\,d\nu\\
 \qquad &=\frac{L_H}{H^{d+1}}\int_X
    \sum_{m\in\mathbb Z}\sum_{n\le H}
       \frac{\theta_H(n)}{L_H}\,
       g_{0,z}(m)\prod_{i=1}^{\ell}g_{i,z}(m+P_i(n))\,d\nu(z).
\end{align*}
Apply \cite[Theorem~4.1]{Teravainen} to the inner sum with weight
$\theta_H/L_H$. The support constant $C$ depends only on the
fixed polynomials $P_i$ and is absorbed into the implied constant.
 By \(L^1\) duality, for some \(K\ge1\) depending
only on \(d\) this gives
\[
 \|\mathcal A_H^F(\mathbf f)-\mu_F\mathcal A_H^1(\mathbf f)\|_1\ll_{\mathbf P}
 L_H\left(H^{-1}+L_H^{-1}\|\theta_H\|_{u^{d+1}[H]}\right)^{1/K}.
\]
Taking \(B>K(A+1)\) yields~\eqref{e:ergodic-L1-log}.

Set \(H_j=\lfloor\exp(j^{1/3})\rfloor\). Taking \(A>3\)
in~\eqref{e:ergodic-L1-log} makes 
\[\sum_j\|\mathcal A_{H_j}^F(\mathbf f)-\mu_F\mathcal A_{H_j}^1(\mathbf f)\|_1 <\infty,\] 
so \(\mathcal A_{H_j}^F(\mathbf f)-\mu_F\mathcal A_{H_j}^1(\mathbf f)\to0\) almost everywhere. For \(H_j\le H\le H_{j+1}\)
the bound \(|\theta_H(n)|\ll_F\log(H+2)\) gives
\begin{align*}
     &|\mathcal A_{H}^F(\mathbf f)(x)-\mu_F\mathcal A_{H}^1(\mathbf f)(x)-\mathcal A_{H_j}^F(\mathbf f)(x)+\mu_F\mathcal A_{H_j}^1(\mathbf f)(x)|\\
 &\ll_F \frac{H_{j+1}-H_j}{H_j}\log(H_{j+1}+2)
 =O_F(j^{-1/3}).
\end{align*}
Here the same summands are used at both scales. This gives
\[\mathcal A_{H}^F(\mathbf f)(x)-\mu_F\mathcal A_{H}^1(\mathbf f)(x)\to0, \qquad \textup{almost everywhere.} \] 
For such an \(x\), write
\[a_p(x)=(F(Z(p-1))-\mu_F)\prod_i f_i(T^{P_i(p-1)}x).\]
The von Mangoldt sum defining \(H(\mathcal A_{H}^F(\mathbf f)(x)-\mu_F\mathcal A_{H}^1(\mathbf f)(x))\), after replacing \(n+1\)
by the prime variable, gives
\[B(H):=\sum_{p\le H}(\log p)a_p(x)=o(H).\] Indeed, the higher
prime powers cost \(O_F(H^{1/2}(\log H)^2)\) and the endpoint
costs \(O_F(\log H)\). Partial summation now gives
\[
 \sum_{p\le H}a_p(x)
 =\frac{B(H)}{\log H}+\int_2^H\frac{B(t)}{t(\log t)^2}\,dt
 =o(H/\log H).
\]
The prime number theorem yields~\eqref{e:smooth-prime-comparison}.
\end{proof}

For the bracket condition, we apply Theorem~\ref{t:pointwise}
on a product system in which two circle rotations supply the
linear and quadratic phases.

\begin{proof}[Proof of Corollary~\ref{c:bracket-pointwise}]
Write $n=p-1$. We first show that for every fixed real polynomial
$\psi$ of degree at most two and every nonzero integer $h$,
\begin{equation}\label{e:quadratic-twist}
 \frac1{\pi(N)}\sum_{p\leq N}
 \e\bigl(-hZ(n)+\psi(n)\bigr)
 \prod_{i=1}^{\ell}f_i(T^{P_i(n)}x)
 \longrightarrow 0
 \qquad\text{for almost every }x.
\end{equation}
As in the proof of Corollary~\ref{c:bracket-recurrence}, adjoin
$n$ and $n^2$ if their degrees are absent, with constant inputs
at the added shifts. Write the linear and quadratic shifts as
$an$ and $bn^2+cn$. Up to a constant, write 
$\psi(n)=\omega_1 an+\omega_2(bn^2+cn)$ for suitable real
$\omega_1,\omega_2$.

Write $Q_1,\ldots,Q_r$ for the enlarged family, with $Q_1(n)=an$
and $Q_2(n)=bn^2+cn$. Let $g_i$ be the original input at $Q_i$,
or $1$ if the shift was added. On $X\times\mathbb T^2$, set
\[
 S(x,u,v)=(Tx,u+\omega_1,v+\omega_2),\quad
 G_1(x,u,v)=g_1(x)\e(u),\quad
 G_2(x,u,v)=g_2(x)\e(v),
\]
and take $G_i(x,u,v)=g_i(x)$ for $i\geq3$. Since
$S^m(x,u,v)=(T^m x,u+m\omega_1,v+m\omega_2)$, we have
\begin{align}\label{e:prodidentity}
 \prod_{i=1}^{r}G_i(S^{Q_i(n)}(x,u,v))
 &=\e\bigl(u+v+\omega_1Q_1(n)+\omega_2Q_2(n)\bigr)
   \prod_{i=1}^{\ell}f_i(T^{P_i(n)}x)\\
 &=\e\bigl(u+v-\psi(0)\bigr)\e(\psi(n))
   \prod_{i=1}^{\ell}f_i(T^{P_i(n)}x)\nonumber.
\end{align}

For a fixed open interval $K\subset\mathbb T$, apply
Theorem~\ref{t:pointwise} to this product system. The restricted
and unrestricted averages have the same limit, and
$\pi_{Z,K}(N)/\pi(N)\to |K|$. Multiplying the restricted average
by this ratio and subtracting $|K|$ times the unrestricted
average gives
\[
 \frac1{\pi(N)}\sum_{p\leq N}
 \bigl(\mathbf1_K(Z(n))-|K|\bigr)
 \prod_{i=1}^{r}G_i(S^{Q_i(n)}(x,u,v))
 \longrightarrow 0
\]
for almost every $(x,u,v)$. By \eqref{e:prodidentity} and
Fubini's theorem, we may remove the factor
$\e(u+v-\psi(0))$ and obtain the same convergence for almost
every $x$:
\begin{align}\label{diffsumtozero}
    \frac1{\pi(N)}\sum_{p\leq N}
 \bigl(\mathbf1_K(Z(n))-|K|\bigr)\e(\psi(n))
 \prod_{i=1}^{\ell}f_i(T^{P_i(n)}x)
 \longrightarrow 0 .
\end{align}
For any fixed $t_0\in\mathbb T$ and $0<\delta<1/2$, the scalar
asymptotic for the open interval of radius $\delta$ about $t_0$
gives
\begin{align*}
   & \limsup_{N\to\infty}
 \frac{|\{p\leq N: Z(p-1)\equiv t_0\pmod 1\}|}{\pi(N)}\\
 &\leq
 \lim_{N\to\infty}
 \frac{|\{p\leq N:\|Z(p-1)-t_0\|_{\mathbb T}<\delta\}|}{\pi(N)}
 =2\delta,
\end{align*}
where both counts are over primes. Letting $\delta\to0$ shows
that adding an endpoint to an interval changes its weighted
average by $o(1)$, since the summands are bounded. The centered
estimate therefore holds for half-open intervals as well.

For an integer $m>|h|$, approximate $\e(-ht)$ on each interval
$[j/m,(j+1)/m)$ by $\e(-hj/m)$, with uniform error at most
$2\pi|h|/m$. Since $\sum_{j=0}^{m-1}\e(-hj/m)=0$, we have
\[
 \sum_{j=0}^{m-1}\e(-hj/m)\mathbf1_{[j/m,(j+1)/m)}(t)
 =
 \sum_{j=0}^{m-1}\e(-hj/m)
 \left(\mathbf1_{[j/m,(j+1)/m)}(t)-\frac1m\right).
\]
By the centered interval estimate and linearity, the average
with this step function in place of $\e(-hZ(n))$ tends to zero
for every $m>|h|$ on a common set of full measure. Combining \eqref{diffsumtozero} with the uniform approximation
bound gives
\[
 \limsup_{N\to\infty}
 \left|
 \frac1{\pi(N)}\sum_{p\leq N}
 \e\bigl(-hZ(n)+\psi(n)\bigr)
 \prod_{i=1}^{\ell}f_i(T^{P_i(n)}x)
 \right|
 \leq \frac{2\pi|h|}{m}\prod_{i=1}^{\ell}\|f_i\|_\infty.
\]
Letting $m\to\infty$ gives \eqref{e:quadratic-twist}.

Now put $t=\{\beta n\}$ and $\theta=\alpha\beta/2$. Since
$\lfloor\beta(n+1)\rfloor=\lfloor\beta n\rfloor+
\mathbf1_{[1-\beta,1)}(t)$, we have
\[
 Z_{\mathrm{br}}(n)+Z(n)
 =\theta(n^2+n)-\alpha t
  +\alpha(n+1)\mathbf1_{[1-\beta,1)}(t).
\]
For $t=\{\beta n\}$, exponentiating the identity above gives
\begin{align*}
 \e(hZ_{\mathrm{br}}(n))
 &=\e\bigl(-hZ(n)+h\theta(n^2+n)\bigr)\\
 &\quad\times\bigl(
 \mathbf1_{[0,1-\beta)}(t)\e(-h\alpha t)
 +\e(h\alpha n)\mathbf1_{[1-\beta,1)}(t)\e(h\alpha(1-t))
 \bigr).
\end{align*}
The two coefficients in parentheses, viewed as functions of $t$,
are continuous away from $0$ and $1-\beta$. Replace each by a
continuous function of modulus at most one that agrees with it
outside the $\delta$-neighborhoods of those two points.
Equidistribution of $\{\beta(p-1)\}$ along the primes gives
\[
 \lim_{N\to\infty}\frac1{\pi(N)}
 \#\{p\leq N:\ p\in\mathbb P,\
       \min(\|t\|_{\mathbb T},\|t-(1-\beta)\|_{\mathbb T})<\delta\}
 \leq 4\delta.
\]
Since each coefficient changes by at most two, the limit superior
of the absolute error in the weighted average is at most
$16\delta\prod_i\|f_i\|_\infty$.

Now approximate the continuous coefficients by trigonometric
polynomials. Each mode $\e(kt)=\e(k\beta n)$ gives a test in
\eqref{e:quadratic-twist} with
$\psi(n)=h\theta(n^2+n)+k\beta n$ or
$\psi(n)=h\theta(n^2+n)+(h\alpha+k\beta)n$.
The corresponding averages tend to zero. Passing to the uniform
limit and then letting $\delta\to0$ proves the cancellation. Thus \eqref{e:quadratic-twist} gives
\begin{align}\label{e:limitofprimesexp}
    \frac1{\pi(N)}\sum_{p\leq N}\e(hZ_{\mathrm{br}}(p-1))
 \prod_{i=1}^{\ell}f_i(T^{P_i(p-1)}x)
 \longrightarrow 0
 \qquad(h\ne0)
\end{align}
on a common set of full measure, using countable sequences of
approximations.

For a trigonometric polynomial $F$, the constant Fourier mode
contributes $(\int_{\mathbb T}F)A_N^T(\mathbf f)(x)$, and every
nonzero mode tends to zero by \eqref{e:limitofprimesexp}. Hence
\[
 \frac1{\pi(N)}\sum_{p\leq N}F(Z_{\mathrm{br}}(p-1))
 \prod_{i=1}^{\ell}f_i(T^{P_i(p-1)}x)
 =\left(\int_{\mathbb T}F\right)A_N^T(\mathbf f)(x)+o(1).
\]
This extends to each fixed continuous $F$ by uniform approximation.
Indeed, replacing $F$ by a trigonometric polynomial changes each
side by at most the uniform approximation error times
$\prod_i\|f_i\|_\infty$.

Use the Lipschitz bounds $F_j^-\leq\mathbf1_J\leq F_j^+$ from the
proof of Theorem~\ref{t:pointwise}. The continuous comparison for
$F_j^+-F_j^-$ with all inputs equal to one bounds the limit superior
of the error from replacing $F_j^-$ by $\mathbf1_J$ by
$2^{-j}\prod_i\|f_i\|_\infty$. Since $\int_{\mathbb T}F_j^-\to |J|$,
letting $j\to\infty$ on a common set of full measure gives
\[
 \frac1{\pi(N)}
 \sum_{\substack{p\leq N\\p\in\mathcal P_{Z_{\mathrm{br}},J}}}
 \prod_{i=1}^{\ell}f_i(T^{P_i(p-1)}x)
 = |J|A_N^T(\mathbf f)(x)+o(1).
\]
Taking all inputs equal to one gives
$\pi_{Z_{\mathrm{br}},J}(N)/\pi(N)\to |J|$. Dividing the displayed
comparison by this ratio proves the corollary.
\end{proof}

\appendix

\section{The cost of dilation to fixed polynomial shifts}
\label{s:why-Cq}

We construct an example with shifts $y,qy^2$ in the smooth-modulus
range of the density increment. The weight is
$G_q(y)=F(Z(qy))-\mu_F$, with the same $\alpha,\beta$ as in the
main theorems, and the scales satisfy
$$\log q\asymp(\log\log M)^4.$$
At every fixed order $S\geq3$, the full Gowers norm of $G_q$
stays bounded below. Thus the estimate from
\cite[Proposition~3.3]{KMTT}, at an order $S\geq3$ chosen using
Gowers monotonicity if necessary, gives no decay.
The quadratic-phase norm has arbitrary logarithmic decay, yet
a single residue class gives a correlation large enough to make
the bound from dilation to the fixed shifts $n,n^2$ vacuous
when applying \cite[Theorem~4.1]{Teravainen}. On the same scales,
our local generalized von Neumann theorem gives
$T_{N,M,G_q}(\mathbf f)=o(1)$ uniformly over the $1$-bounded
inputs.

\medskip
\noindent\emph{The bound obtained by dilation.}
Let \(N=qM^2\) and let the \(f_i\) be \(1\)-bounded and supported
on \([-N,N]\). To apply the fixed-family theorem with shifts
\(n,n^2\), put \(\widetilde G(qy)=G_q(y)\) and
\(\widetilde f_i(qx)=f_i(x)\), extending both by zero off
\(q\mathbb Z\). With parameter length \(qM\) and spatial scale \((qM)^2=qN\),
the count with shifts \(n,n^2\) is
\[
 \widetilde T
 :=\frac{1}{(qM)^3}
 \sum_{m\in\mathbb Z}\sum_{n=1}^{qM}
 \widetilde G(n)\,\widetilde f_0(m)\,
 \widetilde f_1(m+n)\,\widetilde f_2(m+n^2).
\]
Substituting \(m=qx\) and
\(n=qy\) gives \(m+n=q(x+y)\) and \(m+n^2=q(x+qy^2)\).
Thus the two counts have the same numerator, while their
normalizations are \((qN)(qM)\) and \(NM\), respectively. Hence
\begin{equation}\label{e:lift-form}
 T_{N,M,G_q}(\mathbf f)=q^2\widetilde T,
 \qquad
 \|\widetilde G\|_{u^3[qM]}=q^{-1}\|G_q\|_{u^3[M]}.
\end{equation}
For the second identity, a quadratic test \(\psi(n)\) gives
\(q^{-1}\mathbb E_{y\leq M}G_q(y)\e(-\psi(qy))\) after
substitution, and every quadratic \(\phi(y)\) arises by taking
\(\psi(n)=\phi(n/q)\).

After a fixed normalization of the weight, which costs only
a constant depending on \(F\), \cite[Theorem~4.1]{Teravainen}
therefore gives
\begin{equation}\label{e:lift-teravainen-bound}
 |T_{N,M,G_q}(\mathbf f)|
 \ll_F q^2\left((qM)^{-1}
          +q^{-1}\|G_q\|_{u^3[M]}\right)^{1/K}
\end{equation}
for some fixed \(K\geq1\). In particular, if
\(\|G_q\|_{u^3[M]}\gg_F r_j^{-1}\) and
\(\log r_j=o(\log q)\), the expression on the right is at least
a constant times \(q^{2-1/K}r_j^{-1/K}\), which tends to infinity.
We now construct scales with this lower bound and verify that
the local estimate gives decay on the same scales.

\begin{proposition}\label{p:separating-family}
Let $F:\mathbb T\to\mathbb C$ be Lipschitz and nonconstant.
There are sequences $q=q_j$, $M=M_j$, and $N=qM^2$ in the
smooth-modulus range of Theorem~\ref{t:relative-selector-comparison},
with $\log q\asymp(\log r_j)^2$, such that
\begin{equation}\label{e:separating-global-u}
 \|G_q\|_{u^3[M]}\gg_F r_j^{-1}.
\end{equation}
For every fixed $A,B>0$, we have
$\|G_q\|_{u^3[P]}\ll_{A,B,F}(\log M)^{-A}$ on every progression
$P=I\cap(u+q^b\mathbb Z)$, where $I\subset[M]$ is an integer
interval, $b,u\in\mathbb Z$, $0\leq b\leq B$, and
$|P|\geq(qM)^{1/2}$. Also, $\|G_q\|_{U^S[M]}\gg_{F,S}1$
for each fixed integer $S\geq3$.
\end{proposition}

\begin{proof}
Take \(w=\lfloor\log r_j\rfloor\), let \(q\) be the largest
power of \(\mathcal L(w)\) with \(\log q\leq w^2\), and set
\(M=r_j\lfloor r_{j+1}^{1/4}/r_j\rfloor\). Since
\(\log\mathcal L(w)\leq\log(w!)=O(w\log w)\), maximality of
the power gives
\[
 w^2-\log\mathcal L(w)<\log q\leq w^2.
\]
In particular, \(\log q\sim w^2\). The modulus \(q\) is
\(w\)-smooth and is divisible by \(8\prod_{p\leq w}p\) for
large \(j\), since \(\mathcal L(w)\mid q\). By
\eqref{e:horizontal-approximants}, \(M\sim r_{j+1}^{1/4}\) and
\[
 \log M\sim \tfrac14\exp(\sqrt{\log r_j}),
 \qquad \log q\asymp(\log\log M)^4.
\]
Thus \(w=o((\log N)^c)\) for every fixed \(c>0\), and
\(\log q=o(\sqrt{\log N})\). Together with
\(M=N^{1/2}q^{-1/2}\), these give the required smoothness and
scale conditions.

\medskip
\noindent\emph{The lower bound on \([M]\).}
Fix \(h\ne0\) with \(\widehat F(h)\ne0\) and put \(H=M/r_j\).
To prove \eqref{e:separating-global-u}, it suffices to find a
quadratic polynomial \(Q\) such that
\[
 \mathbb E_{t\leq H}G_q(r_jt)\e(-hQ(r_jt))
 =\widehat F(h)+o(1).
\]
Indeed, since \(r_j\mid M\), expanding \(\mathbf1_{r_j\mid y}\)
into linear characters gives
\[
\begin{aligned}
 \frac{|\widehat F(h)+o(1)|}{r_j}
 &=\left|\mathbb E_{y\leq M}
    \mathbf1_{r_j\mid y}G_q(y)\e(-hQ(y))\right|\\
 &\leq \frac1{r_j}\sum_{a=0}^{r_j-1}
    \left|\mathbb E_{y\leq M}
     G_q(y)\e\bigl(-hQ(y)+ay/r_j\bigr)\right|\leq \|G_q\|_{u^3[M]},
\end{aligned}
\]
where the last inequality holds because each test phase is quadratic.

For \(r_j\mid y\), the number \(a_jqy/r_j\) is an integer, and
\eqref{e:horizontal-approximants} gives
\[0<(\beta-a_j/r_j)qy\ll qM/r_{j+1}=o(1).\] Hence
\(\lfloor\beta qy\rfloor=a_jqy/r_j\) throughout this class.
Substitution into \eqref{e:heisenberg-phase} gives
\begin{equation}\label{e:uniappzq}
 Z(qy)=Q(y)+O\!\left(\frac{(qM)^2}{r_{j+1}}\right),
 \qquad
 Q(y)=-\frac{\alpha a_j}{2r_j}(q^2y^2+qy).
\end{equation}
Since $qM=r_{j+1}^{1/4+o(1)}$, the error in
\eqref{e:uniappzq} is $o(1)$ uniformly on the multiples of $r_j$.
 By the Lipschitz continuity of \(F\), we obtain
\[
 \mathbb E_{t\leq H}G_q(r_jt)\e(-hQ(r_jt))
 =
 \mathbb E_{t\leq H}
   \bigl(F(Q(r_jt))-\mu_F\bigr)\e(-hQ(r_jt))
 +O_F\!\left(\frac{(qM)^2}{r_{j+1}}\right).
\]
It remains to show that \(Q(r_jt)\), \(1\leq t\leq H\), is
equidistributed modulo one as \(j\to\infty\). The continuous test
function \(u\mapsto(F(u)-\mu_F)\e(-hu)\) then gives
\[
 \mathbb E_{t\leq H}
   \bigl(F(Q(r_jt))-\mu_F\bigr)\e(-hQ(r_jt))
 =\int_{\mathbb T}(F(u)-\mu_F)\e(-hu)\,du+o(1)
 =\widehat F(h)+o(1).
\]

To prove equidistribution, fix an integer \(k\ne0\). Since
\(Q(r_jt)=-\frac{\alpha a_j}{2}(q^2r_jt^2+qt)\), we have
\[
 k\bigl(Q(r_j(t+v))-Q(r_jt)\bigr)
 =-k\alpha a_jq^2r_jv\,t
  -\frac{k\alpha a_j}{2}(q^2r_jv^2+qv).
\]
The constant term contributes a factor of modulus one. For large \(j\) we have \(0<a_j<r_j\). For
\(1\leq v<H\), the geometric-series formula and bad approximability
give
\[
\begin{aligned}
 &\left|\frac1H\sum_{t=1}^{H-v}
    \e\bigl(k(Q(r_j(t+v))-Q(r_jt))\bigr)\right|\ll
 \frac{1}{H\|k\alpha a_jq^2r_jv\|_{\mathbb T}}
 \leq \frac{|k|a_jq^2r_jv}{c_\alpha H}.
\end{aligned}
\]
Van der Corput's inequality with shifts up to
\(\lfloor H^{1/3}\rfloor\) therefore gives
\[
 \left|\mathbb E_{t\leq H}\e(kQ(r_jt))\right|^2
 \ll_\alpha H^{-1/3}+|k|a_jq^2r_jH^{-2/3}=o(1),
\]
since \(a_j<r_j\) and \(a_jq^2r_j=H^{o(1)}\).
Weyl's criterion proves equidistribution and hence the claimed
correlation on the multiples of \(r_j\).



\medskip
\noindent\emph{The upper bounds on the tested progressions.}
Set \(X=qM\). Since \(X=r_{j+1}^{1/4+o(1)}\) and
\(\log r_j=o(\log r_{j+1})\), we have
\(r_j\leq X^{1/100}<r_{j+1}\). On a progression \(P\) in the
statement, the arguments \(qy\) of \(Z\) form a progression
of length \(|P|\geq X^{1/2}\) and step \(q^{b+1}\).
This step is at most \(X^{1/10000}\) for large \(j\), since
\(q=X^{o(1)}\) and \(b\leq B\). It is coprime to \(r_j\),
since \(q\) is \(w\)-smooth and \(r_j>w\). A quadratic test
in \(y\) is still quadratic after parametrizing \(P\).
Thus \eqref{e:heisenberg-centered} gives
\[
 \|G_q\|_{u^3[P]}\ll_{A,B,F}(\log X)^{-A}
 \ll_{A,B,F}(\log M)^{-A}.
\]
This estimate controls ordinary quadratic phases.
A localized vertical mode gives a nonzero correlation and
therefore a lower bound for the full Gowers norm.

\medskip

\noindent\emph{The lower bound for full Gowers norms.}
Fix \(h\in\mathbb Z\setminus\{0\}\) with \(\widehat F(h)\ne0\).
Choose a smooth nonnegative \(\chi\) with positive integral,
supported on an interval of length less than \(1/2\) inside \((0,1)\). Apply \eqref{e:heisenberg-equidistribution} at scale \(X=qM\),
with \(a=0\), \(v=q\), and \(K=M\), to the Lipschitz function
\[
 \Phi((x,u,z)\Gamma)
 =\chi(\{u\})\bigl(F(z-x\lfloor u\rfloor)-\mu_F\bigr)
   \e\bigl(-h(z-x\lfloor u\rfloor)\bigr).
\]
On \(g_0^{qy}\Gamma\), the second horizontal coordinate is
\(\{\beta qy\}\) and the vertical coordinate is \(Z(qy)\bmod 1\).
The estimate therefore gives
\begin{align}\label{e:estimateexpsum}
 \mathbb E_{y\leq M}G_q(y)\chi(\{\beta qy\})\e(-hZ(qy))
 &=\int_{H/\Gamma}\Phi\,dm_H+o(1)\\
 &=\left(\int_{\mathbb T}\chi(u)\,du\right)
   \left(\int_{\mathbb T}(F(z)-\mu_F)\e(-hz)\,dz\right)+o(1)\nonumber\\
 &=\widehat F(h)\int_{\mathbb T}\chi(u)\,du+o(1)\nonumber.
\end{align}
For the second equality, integrate first along the vertical fibre,
where translation by \(x\lfloor u\rfloor\) preserves circle measure.
The last equality uses \(h\ne0\). 
For \(S\geq3\), normalized Gowers monotonicity and
\eqref{e:vertical-mode-multiplier} give
\begin{align}\label{e:moduleinvus}
 \left|\mathbb E_{y\leq M}
   G_q(y)\chi(\{\beta qy\})\e(-hZ(qy))\right|
 &\ll_S
 \bigl\|G_q(y)\chi(\{\beta qy\})\e(-hZ(qy))\bigr\|_{U^S[M]}\\
 &\leq \|\widehat\chi\|_{\ell^1}\|G_q\|_{U^S[M]}\nonumber.
\end{align}
The first inequality uses that the normalized \(U^1[M]\)-norm
is the absolute value of the average. The second applies the
multiplier bound with \(a=0\), \(v=q\), and frequency \(-h\).
For all sufficiently large \(j\), estimates \eqref{e:estimateexpsum} and \eqref{e:moduleinvus} imply that 
\[
 \|G_q\|_{U^S[M]}
 \gg_S
 \frac{|\widehat F(h)|\int_{\mathbb T}\chi}
      {\|\widehat\chi\|_{\ell^1}}.
\]
The right-hand side is positive and independent of \(j\), since
\(h\) and \(\chi\) are fixed. This gives the claimed lower bound
for each fixed \(S\geq3\).
\end{proof}

\medskip
\noindent\emph{Comparison of the bounds.}
Taking \(P=[M]\) in Proposition~\ref{p:separating-family}
gives arbitrary logarithmic saving for \(\|G_q\|_{u^3[M]}\).
On the other hand, for any fixed \(c,\sigma>0\),
\[
 q^c\|G_q\|_{u^3[M]}^\sigma
 \gg_{F,\sigma} q^c r_j^{-\sigma}\longrightarrow\infty,
\]
because \(\log r_j=o(\log q)\). In
\eqref{e:lift-teravainen-bound} the term involving the norm
has precisely this form, with \(c=2-1/K\) and \(\sigma=1/K\).

One can reduce the normalization loss by repeating the spatial
inputs on all \(q\) residue classes, setting
\(\widetilde f_i(qx+a)=f_i(x)\) for \(0\leq a<q\).
Each class contributes the same numerator, and the support
of each repeated input lies in \([-2qN,2qN]\).
This replaces \(q^2\) by \(q\) in
\eqref{e:lift-teravainen-bound}, with the same weight norm.
The resulting expression is bounded below by a constant
times \(q^{1-1/K}r_j^{-1/K}\), which tends to infinity for
\(K>1\). At \(K=1\), the expression instead reduces to
\(M^{-1}+\|G_q\|_{u^3[M]}=o(1)\).

Theorem~\ref{t:main} applies directly to the shifts \(y,qy^2\).
Take \(\rho=(\log M)^{-1}\) in \eqref{e:main}. Since
\(q=M^{o(1)}\) and \(b\) is bounded, each progression it tests
satisfies
\[
 |I\cap(u+q^b\mathbb Z)|
 \geq \frac{|I|}{q^b}-1
 \geq Mq^{-A_1-b}(\log M)^{-A_2}-1
 \geq(qM)^{1/2}
\]
for large \(j\). Since $\|G_q\|_\infty\leq2\|F\|_\infty$, apply \eqref{e:main}
to $G_q/(1+2\|F\|_\infty)$. By homogeneity and the local phase
bound in Proposition~\ref{p:separating-family}, we obtain
\[
 |T_{N,M,G_q}(\mathbf f)|
 \ll_{A,F} N^{-\eta}+(\log M)^{-1}+(\log M)^{-A\eta}
 \ll_F (\log M)^{-1},
\]
where we take \(A\eta\geq1\). The bound is uniform over the
stated \(1\)-bounded inputs.


The size of the moduli matters here. In the density increment,
the bound \(\log q_i\ll\delta^{-O(1)}w\), together with
\(w\ll\delta^{-E_1}\), gives \(\log q_i\ll\delta^{-D_0}\) for
some fixed \(D_0>0\). At density
\(\delta\geq(\log\log N)^{-c}\), choosing \(cD_0<1\) gives
\[
 \log q_i\ll(\log\log N)^{cD_0}=o(\log\log N).
\]
Since \(\log N_i\geq(\log N)^{3/4}\) throughout the iteration,
these moduli satisfy \(q_i=(\log N_i)^{o(1)}\). In this range
the logarithmic phase savings absorb any fixed power of \(q_i\).
Theorem~\ref{t:recurrence} reaches densities of order
\((\log N)^{-\kappa}\), the same logarithmic scale as the
unweighted result. 
At the polylogarithmic density threshold, the admissible range
includes moduli beyond every fixed power of $\log N_i$.


\begin{thebibliography}{99}

\bibitem{BanksSiegel}
W.~D.~Banks, On the OpenAI proof that Landau--Siegel zeros do not exist,
\href{https://arxiv.org/abs/2610.15483}{arXiv:2610.15483} (2026).

\bibitem{BergelsonLeibman}
V.~Bergelson and A.~Leibman, Polynomial extensions of van der Waerden's and
Szemer\'edi's theorems, \textit{J. Amer. Math. Soc.} \textbf{9} (1996), 725--753.

\bibitem{Euclid}
Euclid, \emph{Elements}, Book~IX, Proposition~20,
online edition by D.~E.~Joyce, Clark University.
\href{https://mathcs.clarku.edu/~djoyce/java/elements/bookIX/propIX20.html}{Link here.}

\bibitem{FHK}
N.~Frantzikinakis, B.~Host and B.~Kra, The polynomial multidimensional
Szemer\'edi theorem along shifted primes, \textit{Israel J. Math.}
\textbf{194} (2013), 331--348.

\bibitem{Furstenberg}
H.~Furstenberg,
Ergodic behavior of diagonal measures and a theorem of Szemer\'edi
on arithmetic progressions,
\emph{J. Anal. Math.} \textbf{31} (1977), 204--256.

\bibitem{GreenSiegel}
B.~Green, On Landau--Siegel zeros, Hexagon preprint
\href{https://hexagonmath.org/2610.00311v2}{hexagon:2610.00311v2} (2026).

\bibitem{GreenTao}
B.~Green and T.~Tao, The primes contain arbitrarily long arithmetic
progressions, \textit{Ann. of Math. (2)} \textbf{167} (2008), 481--547.

\bibitem{GreenTaoNil}
B.~Green and T.~Tao, The quantitative behaviour of polynomial orbits
on nilmanifolds, \emph{Ann. of Math. (2)} \textbf{175} (2012), 465--540.

\bibitem{IwaniecKowalski}
H.~Iwaniec and E.~Kowalski, \textit{Analytic number theory},
AMS Colloquium Publications, vol.~53, American Mathematical Society,
Providence, RI, 2004.

\bibitem{KMTT}
B.~Krause, H.~Mousavi, T.~Tao and J.~Ter\"av\"ainen,
Quantitative bounds for sets lacking polynomial progressions with shifted
prime difference, \href{https://arxiv.org/abs/2608.19525v1}{arXiv:2608.19525v1}.

\bibitem{KMTTPrimes}
B.~Krause, H.~Mousavi, T.~Tao and J.~Ter\"av\"ainen,
Pointwise convergence of bilinear polynomial averages over the primes,
\emph{Ergodic Theory Dynam. Systems} \textbf{45} (2025),
no.~12, 3760--3799.

\bibitem{MirekWanWright}
M.~Mirek, R.~Wan and J.~Wright, Pointwise convergence of polynomial multiple
ergodic averages along the primes,
\href{https://arxiv.org/abs/2505.15549v2}{arXiv:2505.15549v2}.

\bibitem{OpenAISiegel}
OpenAI, Uniform exclusion of Landau--Siegel zeros, OpenAI Math Release
preprint, October~1, 2026,
\href{https://github.com/openai/math/blob/main/preprints/Uniform-exclusion-of-Landau-Siegel-zeros-October-1-2026/README.md}{Link here}.

\bibitem{Peluse}
S.~Peluse, Bounds for sets with no polynomial progressions,
\textit{Forum Math. Pi} \textbf{8} (2020), Paper No.~e16, 55~pp.

\bibitem{PP}
S.~Peluse and S.~Prendiville, A polylogarithmic bound in the nonlinear
Roth theorem, \textit{Int. Math. Res. Not. IMRN} (2022), no.~8,
5658--5684.

\bibitem{Sarkozy}
A.~S\'ark\"ozy, On difference sets of sequences of integers. III,
\textit{Acta Math. Acad. Sci. Hungar.} \textbf{31} (1978), 355--386.

\bibitem{ShaoTeravainen}
X.~Shao and J.~Ter\"av\"ainen, The Bombieri--Vinogradov theorem for
nilsequences, \textit{Discrete Anal.} (2021), Paper No.~21, 55~pp.

\bibitem{ShaoWang}
X.~Shao and M.~Wang, Quantitative bounds in a popular polynomial Szemer\'edi
theorem, \textit{Proc. Roy. Soc. Edinburgh Sect. A}, published online 2025,
\href{https://doi.org/10.1017/prm.2025.10103}{doi:10.1017/prm.2025.10103}.

\bibitem{Szemeredi}
E.~Szemer\'edi, On sets of integers containing no $k$ elements in arithmetic
progression, \textit{Acta Arith.} \textbf{27} (1975), 199--245.

\bibitem{TaoSiegel}
T.~Tao, A sketch of the proof of non-existence of Siegel zeroes,
\textit{What's new}, October~19, 2026,
\href{https://terrytao.wordpress.com/2026/10/19/a-sketch-of-the-proof-of-non-existence-of-siegel-zeroes/}{Link here}.

\bibitem{Teravainen}
J.~Ter\"av\"ainen, Pointwise convergence of ergodic averages with M\"obius
weight, \href{https://arxiv.org/abs/2401.03174v2}{arXiv:2401.03174v2}.

\bibitem{TeravainenWang}
J.~Ter\"av\"ainen and M.~Wang,
On the Green--Tao theorem for sparse sets,
arXiv:2603.09281, 2026.

\bibitem{WooleyZiegler}
T.~D.~Wooley and T.~Ziegler, Multiple recurrence and convergence along the
primes, \textit{Amer. J. Math.} \textbf{134} (2012), 1705--1732.
\end{thebibliography}
\end{document}
